\documentclass[11pt]{amsart}
\usepackage{amssymb,mathtools,amsthm,microtype,mathrsfs,amsmath,geometry,enumitem}
\geometry{margin=1in}
\usepackage{xcolor}
\definecolor{darkblue}{rgb}{0.1,0,0.5}
\usepackage[linktocpage=true,colorlinks,linkcolor=darkblue,citecolor=darkblue,urlcolor=darkblue]{hyperref}
\usepackage[noabbrev,capitalise]{cleveref}

\theoremstyle{plain}
\newtheorem{theorem}{Theorem}[section]
\newtheorem{lemma}[theorem]{Lemma}
\newtheorem{proposition}[theorem]{Proposition}
\newtheorem{corollary}[theorem]{Corollary}
\newtheorem{maintheorem}{Theorem}
\renewcommand{\themaintheorem}{\Alph{maintheorem}}
\theoremstyle{definition}
\newtheorem{definition}[theorem]{Definition}
\theoremstyle{remark}
\newtheorem{remark}[theorem]{Remark}
\crefname{maintheorem}{Theorem}{Theorems}
% Environments that share the theorem counter: tell cleveref the right name.
\AddToHook{env/lemma/begin}{\crefalias{theorem}{lemma}}
\AddToHook{env/proposition/begin}{\crefalias{theorem}{proposition}}
\AddToHook{env/corollary/begin}{\crefalias{theorem}{corollary}}
\AddToHook{env/definition/begin}{\crefalias{theorem}{definition}}
\AddToHook{env/remark/begin}{\crefalias{theorem}{remark}}

% number systems
\newcommand{\RR}{\mathbb{R}}
\newcommand{\CC}{\mathbb{C}}
\newcommand{\NN}{\mathbb{N}}
\newcommand{\ZZ}{\mathbb{Z}}
\newcommand{\QQ}{\mathbb{Q}}
\newcommand{\eps}{\epsilon}
% operators
\DeclareMathOperator{\per}{per}
\DeclareMathOperator{\interior}{int}
\DeclareMathOperator{\Var}{Var}
\DeclareMathOperator{\conv}{conv}
\DeclareMathOperator{\val}{val}
\DeclareMathOperator{\rad}{rad}
\newcommand{\EE}{\mathop{\mathbb{E}}}
% real and complex classes
\newcommand{\cc}[1]{\textsf{\upshape #1}}
\newcommand{\PR}{\cc{P}_{\RR}}
\newcommand{\NPR}{\cc{NP}_{\RR}}
\newcommand{\coNPR}{\cc{coNP}_{\RR}}
\newcommand{\PHR}{\cc{PH}_{\RR}}
\newcommand{\PARR}{\cc{PAR}_{\RR}}
\newcommand{\PATR}{\cc{PAT}_{\RR}}
\newcommand{\EXPR}{\cc{EXP}_{\RR}}
\newcommand{\IPR}{\cc{IP}_{\RR}}
\newcommand{\BIPR}{\cc{BIP}_{\RR}}
\newcommand{\AMR}{\cc{AM}_{\RR}}
\newcommand{\PC}{\cc{P}_{\CC}}
\newcommand{\NPC}{\cc{NP}_{\CC}}
\newcommand{\coNPC}{\cc{coNP}_{\CC}}
\newcommand{\PARC}{\cc{PAR}_{\CC}}
\newcommand{\VP}{\cc{VP}}
\newcommand{\VNP}{\cc{VNP}}
% border complexity and Laurent series
\newcommand{\Leps}{\RR((\eps))}
\newcommand{\bsize}{\underline{S}}
% the truncated power series ring of the Newton lemma
\newcommand{\Rtrunc}{\mathcal{R}}
% facts (S1)-(S8)
\newcommand{\Sfact}[1]{\textup{(S#1)}}


\title{$\NPR\neq\PARR$, $\Sigma_2^{\RR}\neq\PARR$, and $\BIPR\neq\PARR$}
\author[Samuel Everett]{Samuel Everett\\ \tiny{University of Chicago}}
\date{}

\begin{document}

\begin{abstract}
We prove separations of complexity classes in the Blum--Shub--Smale model of computation over the real and complex numbers. Parallel polynomial time is denoted $\PARR$. The inclusions in \textup{(1)--(4)} are known, and we prove that they are strict. In \textup{(5)} we prove the first equality; the other relations in \textup{(5)} are known.
\begin{enumerate}[label=\textup{(\arabic*)},leftmargin=2.2em,itemsep=1pt,topsep=3pt]
\item $\NPR\subsetneq\PARR$: nondeterministic polynomial time over the reals is a proper subclass of parallel polynomial time.
\item $\coNPR\subsetneq\PARR$ and $\Sigma_2^{\RR}\subsetneq\PARR$, where $\Sigma_2^{\RR}$ is the second existential level of the real polynomial hierarchy.
\item $\NPC\subsetneq\PARC$ and $\coNPC\subsetneq\PARC$ over the complex numbers.
\item $\BIPR\subsetneq\PARR$, where $\BIPR$ is the class of Ivanov and de Rougemont defined by interactive proofs with a real verifier and Boolean messages.
\item $\AMR[O(1)]=\NPR\subsetneq\AMR[\mathrm{poly}]=\IPR$: a constant number of rounds of public coins gives exactly the languages with short real certificates, and polynomially many rounds give a larger class.
\end{enumerate}
In every case the separating language is the graph of the permanent. The language $\{(X,Y):Y=\per_m(X)\}$ lies in parallel polynomial time, outside the classes in \textup{(1)--(4)} and outside $\AMR[O(1)]$. Verifiers for it of the kinds that define $\NPR$, $\coNPR$, $\NPC$, $\coNPC$ and $\BIPR$, uniform or not and with arbitrary constants, need time $2^{\Omega((\log m)^{4/3})}$. For $\Sigma_2^{\RR}$ the same bound holds for the running time plus the total length of the two certificate blocks. The respective bounds for $\NPR$, $\coNPR$ and $\Sigma_2^{\RR}$ also hold for each threshold $\{Y\le\per_m(X)\}$ and $\{Y\ge\per_m(X)\}$. A public-coin protocol for the graph with real messages and a polynomial-time verifier needs $\Omega(\log\log m)$ rounds. The proofs combine the lower bound for the border complexity of the permanent with upper bounds such as the following.
\begin{enumerate}[label=\textup{(\arabic*)},leftmargin=2.2em,itemsep=1pt,topsep=3pt,start=6]
\item Let $f$ be a real polynomial of degree $d$ in $n$ variables. If a real computation tree of depth $t$ verifies its graph with one existential or one universal block of certificates, then the border complexity of $f$ is polynomially bounded in $t,d,n$. With an existential block of length $k$ followed by a universal block of length $l$, the bound is polynomial in $t+k+l,d,n$.
\end{enumerate}
Result \textup{(6)} extends a theorem of B\"urgisser from deterministic verification to these three kinds of nondeterministic verification. The proofs of \textup{(6)} and of the first equality in \textup{(5)} use no lower bound. All results are formalized in Lean~4, with the results cited from the literature as explicit hypotheses.
\end{abstract}
\maketitle

\tableofcontents

\section{Introduction}\label{sec:intro}

\subsection{Computation over the reals and its complexity classes}

In the model of Blum, Shub and Smale \cite{BSS89,BCSS98} a machine computes with real numbers. Its registers hold real numbers, an input is a vector $x\in\RR^n$ of size $n$, and one arithmetic operation $+,-,\times,/$ or one comparison $v\ge0$ costs one step, whatever the numbers involved. The program may contain finitely many real constants. This is the natural model for algorithms of numerical analysis, optimisation and computational geometry, and it carries a complexity theory parallel to the classical one. The classes that occur in this paper are the following; precise definitions are in \cref{sec:prelim}.

\begin{itemize}[leftmargin=1.6em,itemsep=2pt]
\item $\PR$ consists of the languages $L\subseteq\RR^\infty=\bigsqcup_n\RR^n$ decided by a machine in polynomial time.
\item $\NPR$ consists of the languages with short real certificates: $x\in L$ exactly when some vector $z$ of polynomially many \emph{real numbers} makes a polynomial-time machine accept $(x,z)$. Deciding whether a system of real polynomial equations has a solution is complete for $\NPR$ \cite{BSS89}. The class $\coNPR$ consists of the complements.
\item $\Sigma_2^{\RR}$ consists of the languages defined by $\exists z\,\forall u$ over a polynomial-time machine, with $z$ and $u$ real vectors of polynomial length. It is the second existential level of the real polynomial hierarchy $\PHR$.
\item $\PARR$, parallel polynomial time, consists of the languages decided by uniform families of algebraic circuits of polynomial depth and at most exponential size. It is the real counterpart of polynomial space.
\end{itemize}

These classes form a chain
\begin{equation}\label{eq:chain}
\PR\ \subseteq\ \NPR\ \subseteq\ \Sigma_2^{\RR}\ \subseteq\ \PHR\ \subseteq\ \PARR ,
\end{equation}
whose last inclusion is quantifier elimination in parallel polynomial time \cite{Ren92,Cuc93,BCSS98}. Over the complex numbers, where machines test only $v=0$, the classes $\PC$, $\NPC$, $\coNPC$ and $\PARC$ are defined in the same way; Hilbert's Nullstellensatz is complete for $\NPC$, and whether $\PC=\NPC$ is the question of Shub and Smale \cite{SS95,BCSS98}.

Few separations are known. Above $\PARR$, the inclusions $\PARR\subseteq\PATR$ and $\PARR\subseteq\EXPR$ are strict, where $\PATR$ is polynomial alternating time \cite{Cuc93,BCSS98}; see the survey \cite{MM97}. Baartse and Meer proved $\PARR\subsetneq\IPR$ for interactive proofs with real messages \cite{BM16}. Inside the chain \eqref{eq:chain}, however, no inclusion was known to be strict until recently. The question whether $\NPR=\PARR$ is the real analogue of the classical question whether $\cc{NP}=\cc{PSPACE}$.

\subsection{The permanent and border complexity}

A second algebraic theory of hardness, due to Valiant \cite{Val79}, concerns the size of arithmetic circuits for families of polynomials. Its central example is the permanent
\[
\per_m(X)=\sum_{\sigma\in S_m}\ \prod_{i=1}^mX_{i\sigma(i)}
\]
of an $m\times m$ matrix of variables, which is complete for the class $\VNP$. A circuit builds a polynomial from variables and constants by additions and multiplications, and its size is the number of gates. The \emph{border size} of a polynomial $f$ allows the constants of the circuit to depend on a parameter $\eps$ and asks only that the circuit compute $f$ up to terms that vanish with $\eps$ (\cref{def:border}); it is never larger than the circuit size. The following lower bound was proved recently.

\medskip
\noindent\textbf{Theorem} (\cite{Cla26}; \cref{thm:lowerbound} below)\textbf{.} \emph{For all large $m$, the border size of $\per_m$ over $\RR$ is greater than $2^{c(\log m)^{4/3}}$, for an absolute constant $c>0$.}
\medskip

The two theories meet in a theorem of B\"urgisser \cite{Bur04}: if the graph $\{Y=f(X)\}$ of a real polynomial of degree $d$ is decided by a computation tree of depth $t$, then the border size of $f$ is polynomial in $t$ and $d$. In words, for polynomials of moderate degree, \emph{checking} a value is not easier than \emph{computing} it. B\"urgisser deduced that a superpolynomial lower bound for the border size of the permanent implies $\PR\ne\PARR$ \cite{Bur01,Bur04}. With the theorem above this gives $\PR\subsetneq\PARR$, the first strict inclusion inside the chain \eqref{eq:chain}.

A deterministic tree is the weakest kind of verifier. This paper asks what happens when the verifier is helped: by a certificate of real numbers, by a universal quantifier, by an alternation of the two, or by a prover with whom it interacts. Throughout, the test language is
\[
\mathrm{PERM}=\bigl\{(X,Y)\in\RR^{m^2}\times\RR:\ Y=\per_m(X),\ m\ge1\bigr\},
\]
the graph of the permanent. It lies in $\PARR$ by Ryser's formula (\cref{lem:ryser}).

\subsection{Main results}

All lower bounds below are stated for \emph{trees} (\cref{def:tree}), the nonuniform version of a machine: one tree for each input size, with arbitrary real constants. A machine running in time $T$ unrolls into trees of depth $O(T)$, so the bounds hold for machines, with or without uniformity.

\begin{maintheorem}[$\NPR\ne\PARR$]\label{thm:A}
The language $\mathrm{PERM}$ is in $\PARR$ and not in $\NPR$. More precisely, suppose that for every $m$ a tree of depth $t_m$ on inputs $(X,Y,z)$ accepts some real certificate $z$ exactly when $Y=\per_m(X)$. Then
\[
t_m\ \ge\ 2^{\Omega((\log m)^{4/3})} .
\]
The same holds for the threshold languages $\{Y\le\per_m(X)\}$ and $\{Y\ge\per_m(X)\}$. Consequently $\NPR\subsetneq\PARR$.
\end{maintheorem}

\Cref{thm:A} says that the value of a real permanent has no short proof: no list of polynomially many real numbers, however cleverly chosen and however it depends on the matrix, can be checked in polynomial time to certify that $Y=\per_m(X)$. The certificate may contain arbitrary real numbers, and the verifier may use arbitrary real constants and may be different for every $m$. In the classical setting the corresponding statement, that the permanent has no $\cc{NP}$ certificates, is not known; it is only known that such certificates would collapse the polynomial hierarchy, by Toda's theorem \cite{Tod91}. \Cref{thm:A} is proved as \cref{cor:np}.

\begin{maintheorem}[universal certificates and the second level]\label{thm:B}
The language $\mathrm{PERM}$ is not in $\coNPR$ and not in $\Sigma_2^{\RR}$, and the threshold languages $\{Y\le\per_m(X)\}$ and $\{Y\ge\per_m(X)\}$ are not in $\coNPR$. A tree of depth $t_m$ that accepts every $z$ exactly when $Y=\per_m(X)$ has $t_m\ge2^{\Omega((\log m)^{4/3})}$. A tree of depth $t_m$ for which
\[
Y=\per_m(X)\quad\Longleftrightarrow\quad\exists z\ \forall u\ \text{the tree accepts }(X,Y,z,u),
\]
with $z$ and $u$ real vectors of lengths $k_m$ and $l_m$, has $t_m+k_m+l_m\ge2^{\Omega((\log m)^{4/3})}$. Consequently $\coNPR\subsetneq\PARR$ and $\Sigma_2^{\RR}\subsetneq\PARR$.
\end{maintheorem}

The first part says that there is also no short \emph{refutation} of a wrong value of the permanent that could be found for every wrong value. The second part reaches the second level of the hierarchy: one alternation of real quantifiers does not help either. Whether the permanent lies outside the whole hierarchy $\PHR$, which would give $\PHR\subsetneq\PARR$, remains open (\cref{sec:open}). \Cref{thm:B} is proved as \cref{cor:conp,cor:second}.

\begin{maintheorem}[the complex numbers]\label{thm:C}
Over $\CC$, the language $\{(X,Y)\in\CC^{m^2}\times\CC:Y=\per_m(X)\}$ is in $\PARC$ and in neither $\NPC$ nor $\coNPC$; the corresponding trees have depth $2^{\Omega((\log m)^{4/3})}$. Consequently $\NPC\subsetneq\PARC$ and $\coNPC\subsetneq\PARC$.
\end{maintheorem}

Over $\CC$ a certificate is a solution of a system of polynomial equations, so \cref{thm:C} says that no system of polynomially many equations of polynomial complexity in $(X,Y,z)$ is solvable in $z$ exactly when $Y=\per_m(X)$. It is proved as \cref{cor:conp}.

The next two results concern interactive proofs. In an interactive proof a polynomial-time verifier exchanges messages with a computationally unbounded prover and flips coins; an input in the language must be accepted with probability at least $2/3$ for a suitable prover, and an input outside it with probability at most $1/3$ for every prover. Over the reals the verifier is a real machine and the coins are discrete. Ivanov and de Rougemont \cite{IdR99} studied the class $\BIPR$ in which all messages are Boolean. They proved that it equals parallel polynomial time in the additive model, a real version of Shamir's theorem \cite{Sha92}, and that $\BIPR\subseteq\PARR$ in the model with multiplication. Baartse and Meer \cite{BM16} studied the class $\IPR$ in which the messages are real numbers and proved $\PARR\subsetneq\IPR$; in particular the protocol of Lund, Fortnow, Karloff and Nisan \cite{LFKN92} verifies the real permanent with real messages.

\begin{maintheorem}[Boolean messages]\label{thm:D}
Every interactive protocol with a real verifier, discrete coins and Boolean messages that decides $\mathrm{PERM}$ has verifier time $2^{\Omega((\log m)^{4/3})}$, for any number of rounds and with private coins. Consequently $\BIPR\subsetneq\PARR$ in the model with multiplication.
\end{maintheorem}

\Cref{thm:D} shows that the inclusion of Ivanov and de Rougemont is strict as soon as the verifier can multiply, and it explains why the protocol for the permanent must send real numbers: a prover that can only send bits, or more generally numbers from a fixed countable set (\cref{rem:discrete}), cannot convince a real verifier of the value of a permanent. It is proved as \cref{thm:bip}.

For real messages we consider public-coin protocols, in which the verifier Arthur sends only its coin flips and the prover Merlin answers with real vectors. Let $\AMR[c]$ be the class of languages with such a protocol of $c$ rounds, and $\AMR[\mathrm{poly}]$ the class with polynomially many rounds; by \cite{BM16} the latter equals $\IPR$.

\begin{maintheorem}[rounds of real interaction]\label{thm:E}
\begin{enumerate}[label=\textup{(\roman*)}]
\item For every constant $c$, $\AMR[c]=\NPR$. Hence
\[
\AMR[O(1)]\ =\ \NPR\ \subsetneq\ \AMR[\mathrm{poly}]\ =\ \IPR .
\]
\item Every public-coin protocol with real messages and a polynomial-time verifier that decides $\mathrm{PERM}$ has $\Omega(\log\log m)$ rounds. In particular $\mathrm{PERM}\notin\AMR[O(1)]$.
\end{enumerate}
\end{maintheorem}

Part (i) uses no lower bound. It says that a bounded number of rounds of randomness can be removed entirely in favour of a single real certificate, because polynomially many coin strings serve all inputs at once and a real machine can carry them in a constant. The strict inclusion in (i) is the theorem of Baartse and Meer. Classically a constant number of rounds also collapses, to two rounds \cite{BM88}, but neither the collapse to $\cc{NP}$ nor the separation from polynomially many rounds is known. Part (ii) shows that the number of rounds in the protocol for the permanent cannot be made constant. \Cref{thm:E} is proved as \cref{thm:am,thm:amper}.

\Cref{thm:A,thm:B,thm:C} and the second part of \cref{thm:E} are consequences of the lower bound for the permanent and of the following theorem, which holds for every polynomial and uses no lower bound.

\begin{maintheorem}[certificates do not help to check a polynomial]\label{thm:F}
Let $f\in\RR[X_1,\dots,X_n]$ have degree at most $d$, and let $U\subseteq\RR^n$ be open and nonempty. Suppose that for $x\in U$ and $y\in\RR$ the relation $y=f(x)$ holds exactly when one of the following holds, for a tree of depth $t$ with arbitrary real constants and real vectors $z$, $u$ of length at most $t$:
\begin{enumerate}[label=\textup{(\roman*)}]
\item the tree accepts $(x,y,z)$ for some $z$;
\item the tree accepts $(x,y,z)$ for all $z$;
\item there is $z$ such that the tree accepts $(x,y,z,u)$ for all $u$.
\end{enumerate}
Then the border size of $f$ is bounded by a polynomial in $t$, $d$ and $n$. In cases \textup{(i)} and \textup{(ii)} the same holds with $y\le f(x)$ or $y\ge f(x)$ in place of $y=f(x)$.
\end{maintheorem}

The case without certificates is B\"urgisser's theorem. \Cref{thm:F} says that his principle, that a polynomial of moderate degree is not easier to check than to compute, survives when the checker is given real certificates, universally quantified real certificates, or both in this order. Case (i) is \cref{thm:graph,thm:threshold}, case (ii) is \cref{thm:onesided}, and case (iii) is \cref{thm:second}. Over the complex numbers cases (i) and (ii) hold for polynomials with real coefficients (\cref{thm:complex}).

In summary, the known relations among these classes are now
\[
\PR\ \subseteq\ \NPR=\AMR[O(1)]\ \subsetneq\ \PARR\ \subsetneq\ \IPR=\AMR[\mathrm{poly}],\qquad
\coNPR,\ \Sigma_2^{\RR},\ \BIPR\ \subsetneq\ \PARR .
\]

\subsection{Ideas of the proofs}

\emph{Existential certificates.} With a certificate the generic-path argument of B\"urgisser breaks down: the polynomials that the tree tests involve the certificate, and no tested polynomial needs to vanish on the graph. After fixing one accepting path and its active tests, the accepted points near a certificate are the common zeros of polynomials $p_1,\dots,p_r$ with small programs, and the certificate of $(x,f(x))$ can be made isolated by freezing some of its coordinates (\cref{lem:isolate}). Then the pair consisting of $f(x)$ and the certificate is the unique zero of $Q=\sum p_i^2$ over $x$. If that zero were nondegenerate, Newton iteration in a truncated power series ring would compute $f$ by a small circuit. In general it is degenerate to arbitrary order. We replace it by the minimiser of $Q+\eps^2|w-c|^2$ for a generic centre $c$, which is a nondegenerate critical point for almost all parameters, and let $\eps\to0$. The limit is taken inside the constants of the circuit, which is exactly what border complexity permits. The step that needs care is that Taylor coefficients, and not only values, converge; this holds on some open ball for every semialgebraic family (\cref{lem:derivatives}). The combination is the local theorem (\cref{thm:local}).

\emph{Universal certificates.} If every $z$ is accepted exactly on the graph, the rejecting certificates live on the two sides of the graph. On one side, a single cell of rejecting computations covers a one-sided neighbourhood of the graph over a ball. We measure how far a certificate stays inside its cell and away from infinity by a margin, and the largest margin $v(x,y)$ is the optimal value of a polynomial program. The local theorem shows that such value functions have Taylor polynomials of small border size at generic points (\cref{prop:value}). The margin tends to zero at the graph, its level sets $\{v=\delta\}$ are graphs of functions $f_\delta\to f$, and the derivatives converge again.

\emph{An existential and a universal block.} Here the good certificates $z$ form, after isolation, a thin graph $\Gamma$ in $(x,y,z)$-space, and every other nearby point has a rejecting certificate. The largest rejecting margin $v$ vanishes exactly on $\Gamma$, but it is a maximum, it may have kinks along $\Gamma$, and it may jump on thin sets, so that its small values need not concentrate near $\Gamma$. We pass to the envelope
\[
e(x;c)=\max_\mu\Bigl[\int v^N\,d\mu-a\Var(\mu)-2a\,|c-\operatorname{mean}(\mu)|^2\Bigr]
\]
over probability measures $\mu$ with finitely many atoms, a regularisation of Lasry--Lions type \cite{LL86}. It is again the largest of finitely many values of polynomial programs, it is continuously differentiable in $c$, and for a suitable power $N$ its zero set is exactly $\Gamma$. A generic linear tilt then produces a nondegenerate critical point in the region where $e$ is smooth.

\emph{Interaction.} Boolean messages cannot depend continuously on the input, so near a generic point of the graph every computation path of the verifier is locally constant unless some tested polynomial vanishes on the graph; freezing the honest prover's strategy then contradicts soundness, and B\"urgisser's theorem on factors applies to the tested polynomial. For a constant number of rounds with real messages, polynomially many coin strings per round suffice for all inputs at once, by a counting bound for sign conditions, and the protocol becomes an existential certificate.

\subsection{What remains open}

The witness $\mathrm{PERM}$ lies outside $\NPR$, so it says nothing about $\PR$ versus $\NPR$. By \cref{thm:F} that question cannot be settled by a lower bound on the border complexity of a polynomial of polynomial degree whose graph is in $\NPR$. \Cref{sec:open} discusses this, the higher levels of the hierarchy, and the relation to a conjecture of B\"urgisser \cite{Bur26}.

\subsection{Organisation}

\Cref{sec:prelim} fixes the model and the tools. \Cref{sec:core} proves the local theorem. \Cref{sec:certificates} treats existential certificates, \cref{sec:universal} universal certificates and the complex field, \cref{sec:second} the second level, and \cref{sec:interaction} interactive proofs.

\section{Preliminaries}\label{sec:prelim}

\subsection{Trees, machines and classes}\label{sec:trees}

We work in the Blum--Shub--Smale model over the ordered field $\RR$ \cite{BSS89,BCSS98}. All lower bounds are proved for the following nonuniform device.

\begin{definition}\label{def:tree}
A \emph{tree} takes an input $w\in\RR^N$. A computation node applies $+$, $-$, $\times$ or $/$ to two operands, each an earlier value, an input coordinate or a real constant. A test node branches on $v\ge 0$ or on $v=0$ for an earlier value $v$. Leaves accept or reject, and an input that reaches a division by zero is rejected. The \emph{depth} is the largest number of nodes on a path from the root to a leaf.
\end{definition}

A machine that runs in time $T$ on inputs of a given size unrolls, for that size, into a tree of depth $O(T)$ whose constants are the constants of the machine. By inserting a test before every division, which at most doubles the depth, we may and do assume that every rejection happens at a leaf.

A language $L\subseteq\RR^\infty=\bigsqcup_n\RR^n$ is in $\PR$ if a polynomial-time machine decides it. It is in $\NPR$ if there are a polynomial-time machine $M$ and a polynomial $q$ such that $x\in L$ exactly when $M$ accepts $(x,z)$ for some $z\in\RR^{q(|x|)}$, and $\coNPR$ consists of the complements. It is in $\Sigma_2^{\RR}$ if $x\in L$ exactly when there is $z$ such that $M$ accepts $(x,z,u)$ for all $u$, with $z$ and $u$ of polynomial length. The class $\PARR$ of parallel polynomial time consists of the languages decided by uniform families of algebraic circuits with sign gates of polynomial depth, and hence of at most exponential size. We use the inclusions
\[
\PR\subseteq\NPR\subseteq\Sigma_2^{\RR}\subseteq\PHR\subseteq\PARR ,
\]
of which the last rests on quantifier elimination in parallel polynomial time for a bounded number of quantifier blocks; see \cite{Ren92,Cuc93,BCSS98}. Among these classes the strict inclusions recorded in the literature are $\PARR\subsetneq\PATR$ and $\PARR\subsetneq\EXPR$ \cite{Cuc93,BCSS98,MM97}, and $\PARR\subsetneq\IPR$ \cite{BM16}. Over $\CC$ the machines branch on $v=0$ only, and $\PC$, $\NPC$, $\coNPC$ and $\PARC$ are defined in the same way.

\subsection{Circuits and border complexity}

A \emph{circuit} over a commutative ring $A$ is a finite list of polynomials over $A$, each of which is a variable, a constant from $A$, or the sum or the product of two earlier entries. Its \emph{size} is the length of the list. The \emph{graph} of a circuit is its list of gate types and wirings, without the values of the constants. A \emph{straight-line program} is a circuit in which we count only the operations $+,-,\times$; a program with $\ell$ operations on $N$ inputs with $p$ constants is a circuit of size at most $3\ell+N+p$.

\begin{definition}\label{def:border}
A circuit over the field $\Leps$ of formal Laurent series \emph{approximates} $f\in\RR[X]$ if it computes a polynomial $f+\eps g$ in which every coefficient of $g$ lies in $\RR[[\eps]]$. The \emph{border size} of $f$ is the least size of a circuit that approximates $f$.
\end{definition}

Our lower bounds all come from the following theorem.

\begin{theorem}[{\cite[Theorems~1.4 and~N.5]{Cla26}}]\label{thm:lowerbound}
There are constants $c>0$ and $m_0$ such that for every $m\ge m_0$ the border size of the permanent $\per_m$ of an $m\times m$ matrix of variables is greater than
\[
m^{c(\log m)^{1/3}}=2^{c(\log m)^{4/3}} .
\]
\end{theorem}

We also use two classical tools.

\begin{theorem}[Baur and Strassen \cite{BS83}]\label{thm:baurstrassen}
If a polynomial has a straight-line program with $\ell$ operations, then the polynomial together with all its first partial derivatives has a straight-line program with $O(\ell)$ operations.
\end{theorem}

\begin{theorem}[B\"urgisser {\cite[Theorem~1.3 and Remark~1.4]{Bur04}}]\label{thm:burgisser}
There is an absolute constant $C$ with the following property. Let $f\in\RR[X_1,\dots,X_n]$ have degree at most $d\ge1$, and suppose that some nonzero multiple of $Y-f(X)$ in $\RR[X,Y]$ has a straight-line program with $\ell$ operations. Then the border size of $f$ is at most $Cd^3(\ell+n)$.
\end{theorem}

\subsection{Semialgebraic tools}\label{sec:semialgebraic}

We use the following standard facts about semialgebraic sets and maps; see \cite{BCR98,vdD98}. Dimension means semialgebraic dimension.

\begin{enumerate}[label=\textup{(S\arabic*)},ref=\textup{(S\arabic*)},leftmargin=3em]
\item\label{S1} Sets and maps defined by first-order formulas over $(\RR,+,\times,<)$ are semialgebraic (Tarski--Seidenberg). Dimension is definable in semialgebraic families.
\item\label{S2} A semialgebraic set $E\subseteq\RR^a\times\RR^b$ has a semialgebraic section over its projection to $\RR^a$ (definable choice).
\item\label{S3} A semialgebraic subset of $\RR$ is a finite union of points and intervals. A semialgebraic function $(0,\delta)\to\RR$ has a limit in $[-\infty,\infty]$ at $0^+$, and on some interval $(0,\delta')$ it is continuous and either constant or strictly monotone (monotonicity theorem).
\item\label{S4} \textup{(a)} A semialgebraic subset of $\RR^n$ has dimension $n$ exactly when it has nonempty interior; if its dimension is smaller, so is that of its closure; a finite union has the largest of the dimensions. \textup{(b)} $\dim(\overline A\setminus A)<\dim A$. \textup{(c)} For $A\subseteq\RR^a\times\RR^b$ with fibers $A_c$ over $c\in\RR^b$, $\dim A=\max_e\bigl(e+\dim\{c:\dim A_c=e\}\bigr)$.
\item\label{S5} (Cell decomposition.) A semialgebraic set is a finite disjoint union of Nash manifolds. Given semialgebraic functions on an open set and $p\ge1$, the open set is a finite disjoint union of cells such that the functions are of class $C^p$ on each open cell. A semialgebraic function on an open set is continuous on an open dense semialgebraic subset. A semialgebraic set in $\RR^n\times\RR$ with finite nonempty fibers over an open set $V$ is, over some open ball in $V$, a finite union of graphs of continuous semialgebraic functions $\lambda_1<\dots<\lambda_q$.
\item\label{S6} The critical values of a Nash map between Nash manifolds form a semialgebraic set of dimension smaller than that of the target (Sard).
\item\label{S7} A semialgebraic function $a\colon(0,\delta)\to\RR$ agrees near $0^+$ with a convergent Puiseux series $\sum_{j\ge j_0}a_j\eta^{j/M}$.
\item\label{S8} A semialgebraic map does not increase dimension, and a semialgebraic bijection preserves it. A semialgebraic function on an open set is Nash on an open dense semialgebraic subset.
\end{enumerate}

For \ref{S7}, the function $a$ satisfies $P(\eta,a(\eta))=0$ for a nonzero polynomial $P$ of positive degree in its second variable. By the Newton--Puiseux theorem the roots of $P(\eta,\cdot)$ for small $\eta>0$ are finitely many convergent Puiseux branches. By \ref{S3} the function $a$ is continuous near $0^+$, and two distinct real branches differ for all small $\eta>0$, so $a$ stays on one branch.

\subsection{The polynomials on a path, and cells}\label{sec:cells}

Fix a path $\pi$ from the root of a tree to a leaf. Represent each value computed on $\pi$ as a fraction $P/D$ of polynomials in the input coordinates. Inputs and constants are $P/1$, and
\[
\frac{P_1}{D_1}\pm\frac{P_2}{D_2}=\frac{P_1D_2\pm P_2D_1}{D_1D_2},\qquad
\frac{P_1}{D_1}\cdot\frac{P_2}{D_2}=\frac{P_1P_2}{D_1D_2},\qquad
\frac{P_1}{D_1}\Big/\frac{P_2}{D_2}=\frac{P_1D_2}{D_1P_2}.
\]
For a tested value put $S=PD$. An input $w$ follows $\pi$ exactly when: \textup{(a)} $P_2(w)\ne0$ at every division on $\pi$; \textup{(b)} at every test $v\ge0$ on $\pi$ we have $S(w)\ge0$ if $\pi$ continues with ``yes'' and $S(w)<0$ if it continues with ``no''; \textup{(c)} at every test $v=0$ we have $P(w)=0$ for ``yes'' and $P(w)\ne0$ for ``no''. Indeed, by induction along $\pi$, condition \textup{(a)} makes every denominator nonzero at $w$, the value is then $P(w)/D(w)$, it has the sign of $S(w)$, and it vanishes exactly when $P(w)$ does.

\begin{lemma}[cells]\label{lem:cells}
Let a tree of depth $t$ be given and let some of its leaves be designated as \emph{marked}. The set of inputs that reach a marked leaf is a finite union of sets of the form
\[
B=Z\cap O,\qquad Z=\{p_1=\dots=p_r=0\},\qquad O=\{q_1>0,\dots,q_m>0\},
\]
called \emph{cells}, where $r,m\le2t$ and, for each cell, all the polynomials $p_i,q_j$ are computed by one straight-line program with at most $8t$ operations from the input coordinates and the constants on one path. The polynomials of a cell involve at most $2t$ input coordinates.
\end{lemma}

\begin{proof}
For a path $\pi$ to a marked leaf let $I_\ge$ be the set of its tests of type \textup{(b)} answered ``yes''. For $J\subseteq I_\ge$ let $B_{\pi,J}$ be the set of inputs that satisfy \textup{(a)}, \textup{(c)}, the strict conditions in \textup{(b)}, $S_j=0$ for $j\in J$ and $S_j>0$ for $j\in I_\ge\setminus J$. The inputs following $\pi$ are the union of the $B_{\pi,J}$. Write each condition $P\ne0$ as $P^2>0$ and each $S<0$ as $-S>0$. The fractions along $\pi$ cost at most $4$ operations per node and each test at most $3$ more. A path has at most $t$ nodes, each with two operands, so at most $2t$ input coordinates occur.
\end{proof}

\section{The analytic core}\label{sec:core}

This section contains the three tools on which every later argument rests, and the local theorem that combines them. Nothing here mentions certificates.

\subsection{Derivatives of semialgebraic families}

Pointwise convergence does not give convergence of derivatives in general. For semialgebraic families it does, on some open ball.

\begin{lemma}\label{lem:derivatives}
Let $V\subseteq\RR^N$ be open and nonempty, let $g\colon V\times(0,\eps_0)\to\RR$ be semialgebraic, and let $g_0\colon V\to\RR$ be continuous with $g(\xi,\eps)\to g_0(\xi)$ as $\eps\to0$ for every $\xi\in V$. Then there are an open ball $V'\subseteq V$ and $\eps'>0$ such that $g$ is of class $C^1$ on $V'\times(0,\eps')$, $g_0$ is of class $C^1$ on $V'$, and for $i=1,\dots,N$
\[
\sup_{\xi\in V'}\Bigl|\frac{\partial g}{\partial\xi_i}(\xi,\eps)-\frac{\partial g_0}{\partial\xi_i}(\xi)\Bigr|\longrightarrow0\qquad(\eps\to0).
\]
\end{lemma}

\begin{proof}
\emph{A box on which $g$ is smooth.} By \ref{S5}, $V\times(0,\eps_0)$ is a finite union of cells with $g$ of class $C^1$ on the open ones. Let $\Sigma$ be the union of the other cells, so $\dim\Sigma\le N$. Let $V^b$ be the set of $\xi\in V$ whose fiber $\Sigma_\xi\subseteq(0,\eps_0)$ contains an interval $(0,e)$. If $V^b$ had nonempty interior, then by \ref{S4}(c) the set $\Sigma$ would have dimension $N+1$. So $\dim V^b<N$, and there is a ball $V_1\subseteq V\setminus\overline{V^b}$. For $\xi\in V_1$ the fiber $\Sigma_\xi$ is a finite union of points and intervals by \ref{S3} and contains no $(0,e)$, so $e(\xi)=\inf(\Sigma_\xi\cup\{\eps_0\})$ is positive. The function $e$ is semialgebraic, hence continuous on an open dense set by \ref{S5}, hence bounded below by some $\eps_1>0$ on a ball $V_2\subseteq V_1$. The box $V_2\times(0,\eps_1)$ misses $\Sigma$. It is connected, so it lies in one open cell and $g$ is $C^1$ on it.

Fix $i$ and let $\mathcal D=\partial g/\partial\xi_i$ on this box, a semialgebraic function. Let $\theta(v)=v/\sqrt{1+v^2}$, a semialgebraic homeomorphism from $\RR$ onto $(-1,1)$, and put $\tilde{\mathcal D}=\theta\circ\mathcal D$.

\emph{Limit values.} Let $\Gamma\subseteq V_2\times(0,\eps_1)\times(-1,1)$ be the graph of $\tilde{\mathcal D}$ and
\[
\Lambda=\{(\xi,v)\in V_2\times[-1,1]:\ (\xi,0,v)\in\overline\Gamma\}.
\]
A copy of $\Lambda$ lies in $\overline\Gamma\setminus\Gamma$, so $\dim\Lambda\le N$ by \ref{S4}(b). Every fiber $\Lambda_\xi$ is nonempty by compactness of $[-1,1]$. By \ref{S4}(c) the set of $\xi$ with infinite $\Lambda_\xi$ has empty interior. By \ref{S4}(a) and \ref{S5} there is a ball $V_4\subseteq V_2$ over which $\Lambda$ is the union of the graphs of continuous semialgebraic functions $\lambda_1<\dots<\lambda_q$ with $q\ge1$. For $\xi\in V_4$ the limit $\tilde\ell(\xi)=\lim_{\eps\to0}\tilde{\mathcal D}(\xi,\eps)\in[-1,1]$ exists by \ref{S3}. It is a semialgebraic function of $\xi$, and $(\xi,\tilde\ell(\xi))\in\Lambda$.

\emph{One sheet.} Suppose $q\ge2$. The index $j(\xi)$ with $\tilde\ell(\xi)=\lambda_{j(\xi)}(\xi)$ is semialgebraic with finitely many values, so it is constant, equal to $j_0$, on a ball $V_5\subseteq V_4$. Take a neighbouring index $b=j_0\pm1$, say $b=j_0+1$; the other case is symmetric. Let $\beta=(\lambda_{j_0}+\lambda_b)/2$. For $\xi\in V_5$ we have $\tilde{\mathcal D}(\xi,\eps)\to\lambda_{j_0}(\xi)<\beta(\xi)$, so
\[
e'(\xi)=\sup\{e\le\eps_1:\ \tilde{\mathcal D}(\xi,\eps'')<\beta(\xi)\text{ for all }\eps''\in(0,e)\}
\]
is positive. It is semialgebraic, hence continuous on an open dense set, so $e'\ge e_0>0$ on a ball $V_6\subseteq V_5$, and $\tilde{\mathcal D}<\beta$ on $V_6\times(0,e_0)$. Take $\xi\in V_6$. Since $(\xi,\lambda_b(\xi))\in\Lambda$, there are $(\xi_m,\eps_m)\to(\xi,0)$ with $\tilde{\mathcal D}(\xi_m,\eps_m)\to\lambda_b(\xi)$. For large $m$ these values are below $\beta(\xi_m)$, and $\beta(\xi_m)\to\beta(\xi)$. So $\lambda_b(\xi)\le\beta(\xi)<\lambda_b(\xi)$, a contradiction. Hence $q=1$, and $\Lambda$ over $V_4$ is the graph of the continuous function $\tilde\ell$.

\emph{Uniform convergence.} Let $K'\subseteq V_4$ be compact. If $\tilde{\mathcal D}(\cdot,\eps)$ did not tend to $\tilde\ell$ uniformly on $K'$, there would be $\eta>0$, $\eps_m\to0$ and $\xi_m\in K'$ with $|\tilde{\mathcal D}(\xi_m,\eps_m)-\tilde\ell(\xi_m)|\ge\eta$. Pass to a subsequence with $\xi_m\to\xi$ and $\tilde{\mathcal D}(\xi_m,\eps_m)\to v$. Then $(\xi,v)\in\Lambda$, so $v=\tilde\ell(\xi)$, while $\tilde\ell(\xi_m)\to\tilde\ell(\xi)$, a contradiction.

\emph{The limit is finite.} The sets $\{\tilde\ell=1\}$, $\{\tilde\ell=-1\}$ and $\{|\tilde\ell|<1\}$ are semialgebraic and cover $V_4$. Suppose $\tilde\ell=1$ on a ball, and let $K'$ be a compact ball inside it. By uniform convergence, for every $M$ there is $e_M$ with $\mathcal D\ge M$ on $K'\times(0,e_M)$. For $\xi$ and $\xi+he_i$ in $K'$ with $h>0$,
\begin{equation}\label{eq:ftc}
g(\xi+he_i,\eps)-g(\xi,\eps)=\int_0^h\mathcal D(\xi+se_i,\eps)\,ds\ \ge\ Mh .
\end{equation}
Letting $\eps\to0$ gives $g_0(\xi+he_i)-g_0(\xi)\ge Mh$ for every $M$, which is impossible. The case $\tilde\ell=-1$ is the same. So $|\tilde\ell|<1$ on some ball $V_8\subseteq V_4$. Put $\ell=\theta^{-1}\circ\tilde\ell$ there. On a compact ball $K''\subseteq V_8$ the values of $\tilde\ell$ stay in a compact subset of $(-1,1)$, on which $\theta^{-1}$ is uniformly continuous. So $\mathcal D(\cdot,\eps)\to\ell$ uniformly on $K''$.

\emph{Identification.} For $\xi$ and $\xi+he_i$ in $K''$, the identity in \eqref{eq:ftc} and uniform convergence give $g_0(\xi+he_i)-g_0(\xi)=\int_0^h\ell(\xi+se_i)\,ds$. Since $\ell$ is continuous, $\partial g_0/\partial\xi_i$ exists and equals $\ell$ in the interior of $K''$.

Treat $i=1,\dots,N$ in turn, each time inside the ball obtained before. Let $V'$ be the interior of the last compact ball and $\eps'=\eps_1$. A continuous function with continuous partial derivatives is $C^1$.
\end{proof}

\begin{corollary}\label{cor:derivatives}
Let $V\subseteq\RR^N$ be open and nonempty, let $W\colon V\times(0,\eps_0)\to\RR^\kappa$ be semialgebraic, and let $\varphi\colon V\to\RR^\kappa$ be continuous with $W(\xi,\eps)\to\varphi(\xi)$ for every $\xi$. Let $D\ge1$. Then there are an open ball $V^*\subseteq V$ and $\eps^*>0$ such that $W$ is $C^1$ on $V^*\times(0,\eps^*)$ and has continuous $\xi$-derivatives of all orders up to $D$ there, $\varphi$ is of class $C^D$ on $V^*$, and for every multi-index $\alpha$ with $|\alpha|\le D$
\[
\sup_{\xi\in V^*}\bigl|\partial_\xi^\alpha W(\xi,\eps)-\partial^\alpha\varphi(\xi)\bigr|\longrightarrow0\qquad(\eps\to0).
\]
\end{corollary}

\begin{proof}
Apply \cref{lem:derivatives} to the coordinates of $W$ one after another, each time restricting to the box found before. The first derivatives are semialgebraic and converge uniformly to the first derivatives of $\varphi$, which are continuous. Apply the lemma to them, and continue. There are finitely many multi-indices of order at most $D$ and finitely many coordinates, so finitely many applications suffice.
\end{proof}

\subsection{From real approximations to approximating circuits}

\begin{lemma}\label{lem:closure}
Fix finitely many circuit graphs. Let $f\in\RR[X]$ and suppose that for every $\eta>0$ one of the graphs computes, for suitable real values of its constants, a polynomial all of whose coefficients are within $\eta$ of the corresponding coefficients of $f$. Then one of the graphs, with suitable constants in $\Leps$, approximates $f$ in the sense of \cref{def:border}.
\end{lemma}

\begin{proof}
One graph serves for a sequence $\eta\to0$; fix it and let it have $p$ constant gates. For $a\in\RR^p$ let $\Psi_a=\sum_\alpha\psi_\alpha(a)X^\alpha$ be the polynomial it computes, so each $\psi_\alpha$ is a polynomial in $a$ with integer coefficients. Let $\mathcal A$ be the finite set of exponents $\alpha$ for which $\psi_\alpha$ is not identically zero or the coefficient $f_\alpha$ of $f$ is nonzero. The set of $(\eta,a)\in(0,1)\times\RR^p$ with $\max_{\alpha\in\mathcal A}|\psi_\alpha(a)-f_\alpha|<\eta$ is semialgebraic, and it projects onto $(0,1)$ because the condition for one $\eta$ implies it for all larger ones. By \ref{S2} it has a semialgebraic section $a(\eta)$. By \ref{S7} there are $M\ge1$ and Laurent series $A_i(s)=\sum_{j\ge j_0}a_{ij}s^j$, convergent for small $s\ne0$, with $a_i(s^M)=A_i(s)$ for small $s>0$. For each $\alpha\in\mathcal A$ the function $s\mapsto\psi_\alpha(a(s^M))$ is then given by the convergent Laurent series $\psi_\alpha(A(s))$, and it differs from $f_\alpha$ by less than $s^M$. A convergent Laurent series whose values tend to $f_\alpha$ as $s\to0^+$ has no terms of negative order and has constant term $f_\alpha$. So with the constants $A_i(\eps)\in\Leps$ the circuit computes $f+\eps g$ with the coefficients of $g$ in $\RR[[\eps]]$.
\end{proof}

\subsection{Newton iteration as a circuit}

For $D\ge1$ and variables $\chi=(\chi_1,\dots,\chi_N)$ let $\Rtrunc=\RR[\chi][\tau]/(\tau^{D+1})$. An element is stored as its $D+1$ coefficients in $\RR[\chi]$. A sum costs $D+1$ gates, a product with a real constant $D+1$ gates, a difference $2(D+1)$ gates, and a product of two elements at most $(D+1)(D+2)$ gates.

\begin{lemma}\label{lem:newton}
Let $E=(E_1,\dots,E_\kappa)$ be polynomials in $(\xi,w)\in\RR^N\times\RR^\kappa$ computed by a straight-line program with $\ell$ operations. Let $\xi_1\in\RR^N$, $w_1\in\RR^\kappa$ and a real $\kappa\times\kappa$ matrix $A$ be given. In $\Rtrunc^\kappa$ put
\[
\Omega^{(0)}=w_1,\qquad \Omega^{(j+1)}=\Omega^{(j)}-A\,E\bigl(\xi_1+\tau\chi,\ \Omega^{(j)}\bigr),
\]
and let $g_i\in\RR[\xi]$ be the sum of the $D+1$ coefficients of $\Omega^{(D)}_i$, with $\chi=\xi-\xi_1$ substituted.
\begin{enumerate}[label=\textup{(\alph*)}]
\item The $g_i$ are computed by a circuit with at most $C\bigl((D+1)^3\ell+(D+1)^2\kappa^2+N\bigr)$ gates plus the constants of the program, whose graph depends only on the graph of the program and on $N$, $\kappa$ and $D$. Each $g_i$ has degree at most $D$. The coefficients of the $g_i$ are polynomial functions of the coefficients of $E$ and of $\xi_1$, $w_1$ and $A$.
\item Suppose $E(\xi_1,w_1)=0$ and $A$ is the inverse of $\partial_wE(\xi_1,w_1)$. Let $\omega$ be the analytic map near $\xi_1$ with $\omega(\xi_1)=w_1$ and $E(\xi,\omega(\xi))=0$. Then $g_i$ is the Taylor polynomial of order $D$ of $\omega_i$ at $\xi_1$.
\end{enumerate}
\end{lemma}

\begin{proof}
(a) Each of the $D$ steps evaluates the program in $\Rtrunc$, at a cost of $O(\ell(D+1)^2)$ gates, multiplies by the constant matrix $A$ and subtracts, at a cost of $O(\kappa^2(D+1))$ gates. Forming $\chi$ and the final sums costs $O(N+\kappa D)$ gates. The wiring never refers to the values of the constants. For every choice of the constants, each element of $\Rtrunc$ that the circuit forms has, as its coefficient of $\tau^j$, a homogeneous polynomial of degree $j$ in $\chi$: this holds for the inputs $\xi_1+\tau\chi$ and for constants, and sums and products preserve it. So the output has degree at most $D$. Evaluating a program in $\Rtrunc$ evaluates the polynomials it computes, so the output depends on the program only through the coefficients of $E$.

(b) Let $\Omega\in\RR[\chi][[\tau]]^\kappa$ be the Taylor series of $\tau\mapsto\omega(\xi_1+\tau\chi)$. Its coefficient of $\tau^j$ is the homogeneous part of degree $j$ of the Taylor series of $\omega$ at $\xi_1$, and $E(\xi_1+\tau\chi,\Omega)=0$ as an identity of power series. We claim $\Omega^{(j)}\equiv\Omega\pmod{\tau^{j+1}}$. This holds for $j=0$. If $\Omega^{(j)}=\Omega+\tau^{j+1}\Delta$, the Taylor expansion of the polynomial map $E$ in $w$ gives
\[
E(\xi_1+\tau\chi,\Omega^{(j)})=\partial_wE(\xi_1+\tau\chi,\Omega)\,\tau^{j+1}\Delta+O(\tau^{2j+2})\equiv A^{-1}\tau^{j+1}\Delta\pmod{\tau^{j+2}},
\]
since $\partial_wE(\xi_1+\tau\chi,\Omega)\equiv A^{-1}\pmod\tau$. So $\Omega^{(j+1)}\equiv\Omega\pmod{\tau^{j+2}}$. After $D$ steps $\Omega^{(D)}\equiv\Omega\pmod{\tau^{D+1}}$.
\end{proof}

\subsection{The local theorem}

\begin{theorem}[local theorem]\label{thm:local}
Let $K\subseteq\RR^N$ be a closed ball, let $\bar B\subseteq\RR^\kappa$ be the closed ball of radius $r_0$ about $w_0$, and let $\lambda\in\RR^\kappa$. Let $Q\in\RR[\xi,w]$ be nonnegative on $K\times\bar B$ and computed by a straight-line program with $\ell$ operations. Let $\varphi\colon K\to\RR^\kappa$ be continuous with $|\varphi(\xi)-w_0|\le r_0/2$ on $K$, and suppose that for every $\xi\in K$
\begin{equation}\label{eq:unique}
\varphi(\xi)\ \text{is the unique maximiser of }\langle\lambda,w\rangle\text{ on }\{w\in\bar B:\ Q(\xi,w)=0\}.
\end{equation}
Let $D\ge1$. Then there are an open ball $V^*\subseteq K$, a semialgebraic set $E^*\subseteq V^*$ of dimension less than $N$, and a circuit graph with $\kappa$ outputs and at most $C\bigl((D+1)^3\ell+(D+1)^2\kappa^2+N\bigr)$ gates, which depends only on the graph of the program and on $N$, $\kappa$ and $D$, such that $\varphi$ is of class $C^D$ on $V^*$ and, for every $\xi_1\in V^*\setminus E^*$, the Taylor polynomial of order $D$ of $\varphi_i$ at $\xi_1$ is a coefficientwise limit of polynomials of degree at most $D$ computed at the $i$-th output with real constants.
\end{theorem}

For $\lambda=0$ condition \eqref{eq:unique} says that $\varphi(\xi)$ is the only zero of $Q(\xi,\cdot)$ in $\bar B$. The zero may be degenerate to any order, and nothing bounds its multiplicity. The proof regularises it.

For $c\in\RR^\kappa$ and $\eps>0$ put
\[
G_{c,\eps}(\xi,w)=Q(\xi,w)-\eps\langle\lambda,w\rangle+\eps^2|w-c|^2 .
\]

\begin{lemma}\label{lem:minimisers}
Fix $c$. There is $\eps_0>0$ such that for $\xi\in K$ and $0<\eps<\eps_0$ every minimiser $w^*$ of $G_{c,\eps}(\xi,\cdot)$ on $\bar B$ lies in the interior of $\bar B$ and satisfies
\begin{equation}\label{eq:critical}
\nabla_wQ(\xi,w^*)-\eps\lambda+2\eps^2(w^*-c)=0 .
\end{equation}
Moreover $|w^*-\varphi(\xi)|\to0$ as $\eps\to0$, uniformly over $\xi\in K$ and over all minimisers.
\end{lemma}

\begin{proof}
Let $C_0=\sup_{\xi\in K}|\varphi(\xi)-c|^2$. Since $Q(\xi,\varphi(\xi))=0$, a minimiser satisfies $G_{c,\eps}(\xi,w^*)\le-\eps\langle\lambda,\varphi(\xi)\rangle+\eps^2C_0$. As $Q\ge0$ and $|w^*-c|^2\ge0$, this gives
\begin{equation}\label{eq:twobounds}
\langle\lambda,w^*\rangle\ge\langle\lambda,\varphi(\xi)\rangle-\eps C_0,\qquad
Q(\xi,w^*)\le\eps\,\langle\lambda,w^*-\varphi(\xi)\rangle+\eps^2C_0\le\eps\,(2r_0|\lambda|+\eps C_0).
\end{equation}
Suppose the convergence is not uniform. Then there are $\eta>0$, $\eps_j\to0$, $\xi_j\in K$ and minimisers $w_j$ with $|w_j-\varphi(\xi_j)|\ge\eta$. Pass to a subsequence with $\xi_j\to\xi^\circ$ and $w_j\to w^\circ$. By \eqref{eq:twobounds} and the continuity of $\varphi$, $Q(\xi^\circ,w^\circ)=0$ and $\langle\lambda,w^\circ\rangle\ge\langle\lambda,\varphi(\xi^\circ)\rangle$. By \eqref{eq:unique}, $w^\circ=\varphi(\xi^\circ)$, which contradicts $|w^\circ-\varphi(\xi^\circ)|\ge\eta$. For small $\eps$ the minimisers are therefore within $r_0/4$ of $\varphi(\xi)$, hence interior, and the gradient of $G_{c,\eps}(\xi,\cdot)$ vanishes there.
\end{proof}

For $c\in\RR^\kappa$ let $\mathrm{Bad}_c$ be the set of $(\xi,\eps)\in\RR^N\times(0,\infty)$ for which some $w\in\RR^\kappa$ satisfies \eqref{eq:critical} and $\det\bigl(\nabla_w^2Q(\xi,w)+2\eps^2I\bigr)=0$.

\begin{lemma}\label{lem:genericcentre}
There is $c$ with $\dim\mathrm{Bad}_c\le N$.
\end{lemma}

\begin{proof}
On the manifold $\RR^N\times(0,\infty)\times\RR^\kappa$ consider the Nash map
\[
\Phi(\xi,\eps,w)=\Bigl(\xi,\ \eps,\ w+\tfrac1{2\eps^2}\bigl(\nabla_wQ(\xi,w)-\eps\lambda\bigr)\Bigr)
\]
into itself. A triple satisfies \eqref{eq:critical} exactly when $\Phi(\xi,\eps,w)=(\xi,\eps,c)$. The differential of $\Phi$ is block triangular with determinant $(2\eps^2)^{-\kappa}\det(\nabla_w^2Q+2\eps^2I)$. So $(\xi,\eps)\in\mathrm{Bad}_c$ exactly when $(\xi,\eps,c)$ is a critical value of $\Phi$. By \ref{S6} the set of critical values has dimension at most $N+\kappa$, and its fiber over $c$ is $\mathrm{Bad}_c$. If every fiber had dimension $N+1$, then \ref{S4}(c) would give dimension $N+1+\kappa$.
\end{proof}

\begin{proof}[Proof of \cref{thm:local}]
Fix $c$ as in \cref{lem:genericcentre} and $\eps_0$ as in \cref{lem:minimisers}. The set of $(\xi,\eps,w)$ with $\xi\in K$, $0<\eps<\eps_0$ and $w$ a minimiser of $G_{c,\eps}(\xi,\cdot)$ on $\bar B$ is semialgebraic by \ref{S1} and projects onto $K\times(0,\eps_0)$. By \ref{S2} there is a semialgebraic map $W\colon K\times(0,\eps_0)\to\bar B$ with $W(\xi,\eps)$ a minimiser. By \cref{lem:minimisers}, equation \eqref{eq:critical} holds at $w=W(\xi,\eps)$ and $W(\cdot,\eps)\to\varphi$ uniformly on $K$.

By \cref{cor:derivatives}, applied on the interior of $K$, there are an open ball $V^*$ and $\eps^*\in(0,\eps_0)$ such that $W$ is $C^1$ on $V^*\times(0,\eps^*)$, $\varphi$ is $C^D$ on $V^*$, and the $\xi$-derivatives of $W(\cdot,\eps)$ up to order $D$ converge uniformly on $V^*$ to those of $\varphi$.

Let $E^*$ be the set of $\xi_1\in V^*$ whose fiber $\{\eps:(\xi_1,\eps)\in\mathrm{Bad}_c\}$ contains an interval $(0,e)$. If $E^*$ had dimension $N$, then $\mathrm{Bad}_c$ would have dimension $N+1$ by \ref{S4}(c). So $\dim E^*<N$. Let $\xi_1\in V^*\setminus E^*$. By \ref{S3} its fiber is a finite union of points and intervals, so all sufficiently small $\eps$ lie outside it. For such $\eps<\eps^*$ put $w_1=W(\xi_1,\eps)$ and
\[
E(\xi,w)=\nabla_wQ(\xi,w)-\eps\lambda+2\eps^2(w-c).
\]
Then $E(\xi_1,w_1)=0$, and $\partial_wE(\xi_1,w_1)=\nabla_w^2Q(\xi_1,w_1)+2\eps^2I$ is invertible because $(\xi_1,\eps)\notin\mathrm{Bad}_c$. By the analytic implicit function theorem there is an analytic map $\omega$ near $\xi_1$ with $\omega(\xi_1)=w_1$ and $E(\xi,\omega(\xi))=0$, and every continuous solution through $(\xi_1,w_1)$ agrees with $\omega$ near $\xi_1$. The map $W(\cdot,\eps)$ is such a solution. So the Taylor polynomial of order $D$ of $\omega_i$ at $\xi_1$ is that of $W_i(\cdot,\eps)$, and as $\eps\to0$ its coefficients converge to those of the Taylor polynomial of $\varphi_i$.

By \cref{thm:baurstrassen} the map $E$ has a straight-line program with $O(\ell+\kappa)$ operations; the construction of Baur and Strassen transforms the program gate by gate, so the graph of the new program depends only on that of the program for $Q$, and the numbers $\eps$, $c$ and $\lambda$ are constants of it. Moreover $\kappa\le2\ell$: if $Q$ did not involve some coordinate $w_j$, then $\varphi(\xi)+te_j$ would be a zero of $Q(\xi,\cdot)$ in $\bar B$ for small $|t|$, with a value of $\langle\lambda,\cdot\rangle$ at least as large for a suitable sign of $t$, against \eqref{eq:unique}; and a program with $\ell$ operations involves at most $2\ell$ input coordinates. \Cref{lem:newton}, with $A$ the inverse of $\partial_wE(\xi_1,w_1)$, computes the Taylor polynomials of the $\omega_i$ by a circuit of the stated size whose graph does not depend on $\eps$ or on $\xi_1$.
\end{proof}

\section{Existential certificates}\label{sec:certificates}

\subsection{Isolating a certificate}

Certificates need not be unique, and the set of certificates of a point may have positive dimension. The next lemma selects, over a small ball, one certificate that is isolated once some of its coordinates are frozen.

\begin{lemma}[isolating a certificate]\label{lem:isolate}
Let $U_1\subseteq\RR^n$ be an open ball and let $F\subseteq U_1\times\RR^k$ be a semialgebraic set that projects onto $U_1$. Then there are a point $q_0=(x_0,z_0)\in F$, an open neighbourhood $N$ of $q_0$, a set $I$ of coordinates of $\RR^k$, a ball $U_3\subseteq U_1$ about $x_0$ and an analytic map $\zeta\colon U_3\to\RR^k$ such that, for every $x\in U_3$, the point $(x,\zeta(x))$ is the only point $(x,z)$ of $F\cap N$ whose coordinates $z_I$ equal those of $z_0$.
\end{lemma}

\begin{proof}
For an open $U'\subseteq U_1$ write $F_{U'}=F\cap(U'\times\RR^k)$. The dimensions of the sets $F_{U'}$, for open balls $U'\subseteq U_1$, are integers between $n$ and $n+k$. Choose a ball $U_2$ for which $\Delta=\dim F_{U_2}$ is smallest. Then $\dim F_{U'}=\Delta$ for every ball $U'\subseteq U_2$.

Write $F_{U_2}$ as a finite disjoint union of Nash manifolds $C_1,\dots,C_s$ by \ref{S5}. Let $\mathrm{Reg}$ be the union, over the $C_i$ of dimension $\Delta$, of the sets $C_i\setminus\bigcup_{j\ne i}\overline{C_j}$. A point of such a set has a neighbourhood that meets no $C_j$ with $j\ne i$. Hence $\mathrm{Reg}$ is open in $F_{U_2}$ and is a Nash manifold of dimension $\Delta$. The remainder $F_{U_2}\setminus\mathrm{Reg}$ lies in the union of the $C_j$ of dimension below $\Delta$ and of the sets $C_i\cap\overline{C_j}\subseteq\overline{C_j}\setminus C_j$ with $j\ne i$. By \ref{S4}(a) and \ref{S4}(b) its dimension is below $\Delta$.

The projection of $\mathrm{Reg}$ to $\RR^n$ meets every ball $U'\subseteq U_2$, since otherwise $F_{U'}\subseteq F_{U_2}\setminus\mathrm{Reg}$ would have dimension below $\Delta$. A semialgebraic set that is dense in $U_2$ has dimension $n$ by \ref{S4}(a). By \ref{S6} the critical values of the projection $\mathrm{Reg}\to\RR^n$ form a set of dimension below $n$. So the image contains a regular value $x_0$. Take any $q_0=(x_0,z_0)\in\mathrm{Reg}$ and a neighbourhood $N\subseteq U_2\times\RR^k$ of $q_0$ with $F\cap N=\mathrm{Reg}\cap N$ on which the projection is a submersion. Then $F\cap N$ is a manifold of dimension $n+\delta$ with $\delta=\Delta-n\ge0$.

The tangent space of $F\cap N$ at $q_0$ maps onto $\RR^n$. Its intersection $V_0$ with $\{0\}\times\RR^k$ has dimension $\delta$. Choose a set $I$ of $\delta$ coordinates of $\RR^k$ such that $z\mapsto z_I$ is injective on $V_0$; it exists because a matrix of rank $\delta$ has a nonzero $\delta\times\delta$ minor. The map $(x,z)\mapsto(x,z_I)$ from $F\cap N$ to $\RR^{n+\delta}$ then has bijective differential at $q_0$. By the inverse function theorem, after shrinking $N$, it is a diffeomorphism onto an open set, which contains $U_3\times\{(z_0)_I\}$ for a small ball $U_3$ about $x_0$. Let $(x,\zeta(x))$ be the preimage of $(x,(z_0)_I)$.
\end{proof}

In the situation of the lemma we \emph{freeze} the coordinates in $I$: we write $z=(a,u)$ with $a=(z_0)_I$ a vector of constants and $u\in\RR^{k'}$, $k'=k-|I|$, the remaining coordinates, and $\zeta(x)=(a,\upsilon(x))$.

\subsection{The graph theorem}

\begin{theorem}[graph theorem]\label{thm:graph}
Let $f\in\RR[X_1,\dots,X_n]$ have degree at most $d$ with $d\ge1$, and let $U\subseteq\RR^n$ be open and nonempty. Suppose a tree of depth $t$ on inputs $(x,y,z)\in\RR^n\times\RR\times\RR^k$ satisfies, for all $x\in U$ and all $y\in\RR$,
\[
\text{the tree accepts }(x,y,z)\text{ for some }z\in\RR^k\quad\Longleftrightarrow\quad y=f(x).
\]
Then the border size of $f$ is at most $C\bigl((d+1)^3\,t+(d+1)^2(\bar k+2)^2+n\bigr)$, where $\bar k=\min(k,2t)$ and $C$ is an absolute constant.
\end{theorem}

The tree may use any real constants, so the theorem is nonuniform. Boolean certificates are the special case in which the tree checks $z_i(z_i-1)=0$. B\"urgisser's theorem on the decision complexity of graphs \cite[Corollary~1.5]{Bur04} is the case $k=0$, with a sharper bound.

\begin{proof}
\emph{Fixing a cell.} Apply \cref{lem:cells} with the accepting leaves marked. For $x\in U$ the point $(x,f(x))$ has a certificate, so $U$ is the union, over the finitely many cells $B$, of the semialgebraic sets $\{x\in U:(x,f(x),z)\in B\text{ for some }z\}$. By \ref{S4}(a) one of them has nonempty interior. Fix that cell $B=Z\cap O$ with polynomials $p_1,\dots,p_r$ and $q_1,\dots,q_m$. Certificate coordinates that do not occur in these polynomials do not affect membership in $B$. Fix them to $0$ and let $k$ denote the number of the others, at most $2t$. Let $U_1$ be an open ball in the chosen set and
\[
F=\{(x,z)\in U_1\times\RR^k:\ (x,f(x),z)\in B\}.
\]

\emph{Isolating.} Apply \cref{lem:isolate} to $F$ and freeze coordinates. Put $w=(y,u)\in\RR^\kappa$ with $\kappa=1+k'$,
\[
Q(x,w)=\sum_{i=1}^rp_i\bigl(x,y,(a,u)\bigr)^2,\qquad\varphi(x)=\bigl(f(x),\upsilon(x)\bigr),\qquad w_0=\varphi(x_0).
\]
By \cref{lem:cells}, $Q$ has a straight-line program with at most $12t$ operations. The sets $O$ and $N$ are open and contain $(x_0,f(x_0),z_0)$ and $q_0$, and $\varphi$ is continuous. So there are $r_0>0$ and a closed ball $K\subseteq U_3$ about $x_0$ such that $|\varphi(x)-w_0|\le r_0/2$ on $K$ and, for $x\in K$ and $|w-w_0|\le r_0$, the point $(x,y,(a,u))$ lies in $O$ and $(x,(a,u))$ lies in $N$.

\emph{The zero is isolated.} We claim that for $x\in K$ and $|w-w_0|\le r_0$
\[
Q(x,w)=0\quad\Longleftrightarrow\quad w=\varphi(x).
\]
If $w=\varphi(x)$, then $(x,f(x),\zeta(x))\in B\subseteq Z$ and $Q(x,w)=0$. Conversely let $Q(x,w)=0$. Every $p_i$ vanishes at $(x,y,(a,u))$, so this point lies in $Z\cap O=B$ and is accepted. Since $x\in U$, the hypothesis gives $y=f(x)$. Then $(x,(a,u))$ is a point of $F\cap N$ whose frozen coordinates are $a$, so $(a,u)=\zeta(x)$ and $u=\upsilon(x)$.

\emph{Conclusion.} The hypotheses of \cref{thm:local} hold with $\xi=x$ as the parameter, in dimension $n$, with $\lambda=0$ and $D=d$. It gives a ball $V^*$, a point $x_1\in V^*\setminus E^*$, and a circuit graph of the stated size whose first output computes, for suitable real constants, polynomials of degree at most $d$ that converge coefficientwise to the Taylor polynomial of order $d$ of $f$ at $x_1$. Since $\deg f\le d$, that Taylor polynomial is $f$. By \cref{lem:closure} the graph approximates $f$.
\end{proof}

\begin{remark}[several outputs]\label{rem:several}
The same proof applies to a polynomial map $f=(f_1,\dots,f_\mu)\colon\RR^n\to\RR^\mu$ and the language $y=f(x)$ with $y\in\RR^\mu$: the block $y$ is part of $w$, and \cref{thm:local} gives every coordinate of $\varphi$. The bound for each $f_i$ is $C\bigl((d+1)^3t+(d+1)^2(\bar k+\mu+1)^2+n\bigr)$.
\end{remark}

\subsection{Thresholds}

\begin{theorem}[threshold theorem]\label{thm:threshold}
Let $f$, $d$, $U$ and the tree be as in \cref{thm:graph}, but suppose that for all $x\in U$ and $y\in\RR$
\[
\text{the tree accepts }(x,y,z)\text{ for some }z\quad\Longleftrightarrow\quad y\le f(x).
\]
Then the border size of $f$ satisfies the bound of \cref{thm:graph}. The same holds with $y\ge f(x)$ in place of $y\le f(x)$.
\end{theorem}

\begin{proof}
The boundary points $(x,f(x))$ are accepted. Define the cell $B$, the set $F$, the frozen coordinates, $Q$, $\varphi$, $K$ and $r_0$ exactly as in the proof of \cref{thm:graph}. Let $x\in K$ and $|w-w_0|\le r_0$. If $Q(x,w)=0$, the point $(x,y,(a,u))$ lies in $B$ and is accepted, so $y\le f(x)$. If moreover $y=f(x)$, then $(x,(a,u))\in F\cap N$ and $u=\upsilon(x)$. Since $Q(x,\varphi(x))=0$, the point $\varphi(x)$ is the unique maximiser of the coordinate $y$ on the zero set of $Q(x,\cdot)$ in the ball. So \cref{thm:local} applies with $\lambda$ the unit vector of the coordinate $y$, and the proof ends as before. For $y\ge f(x)$ replace $(y,f)$ by $(-y,-f)$.
\end{proof}

\subsection{Value functions}

The next proposition is the form of the local theorem used for universal quantifiers. It concerns the optimal value of a polynomial program and does not require the value to be a polynomial.

\begin{proposition}[value functions]\label{prop:value}
Let $Q\in\RR[\xi,s,\nu]$, with $\xi\in\RR^N$, $s\in\RR$ and $\nu\in\RR^\kappa$, be nonnegative and computed by a straight-line program with $\ell$ operations. Let $A\subseteq\RR^N$ be open and nonempty, and suppose that for every $\xi\in A$ the coordinate $s$ attains a maximum $v(\xi)$ on the set $\{(s,\nu):Q(\xi,s,\nu)=0\}$. Let $D\ge1$. Then there are a dense semialgebraic set $A^\circ\subseteq A$ and finitely many circuit graphs, each with at most $C\bigl((D+1)^3\ell+(D+1)^2(\kappa+1)^2+N\bigr)$ gates, such that for every $\xi_1\in A^\circ$ the function $v$ is of class $C^D$ near $\xi_1$ and its Taylor polynomial of order $D$ at $\xi_1$ is a coefficientwise limit of polynomials of degree at most $D$ computed by one of these graphs with real constants.
\end{proposition}

\begin{proof}
The function $v$ is semialgebraic by \ref{S1}, so by \ref{S5} it is continuous on an open dense subset $A_1$ of $A$. Let $U_1$ be any open ball in $A_1$ and put
\[
F=\{(\xi,\nu)\in U_1\times\RR^\kappa:\ Q(\xi,v(\xi),\nu)=0\},
\]
which projects onto $U_1$ because the maximum is attained. Apply \cref{lem:isolate} to $F$ and freeze coordinates, writing $\nu=(a,u)$. Put $w=(s,u)$, $\tilde Q(\xi,w)=Q(\xi,s,(a,u))$, $\varphi(\xi)=(v(\xi),\upsilon(\xi))$ and $w_0=\varphi(\xi_0)$. The map $\varphi$ is continuous near $\xi_0$. Choose $r_0>0$ and a closed ball $K\subseteq U_3$ about $\xi_0$ such that $|\varphi-w_0|\le r_0/2$ on $K$ and $(\xi,(a,u))\in N$ whenever $\xi\in K$ and $|w-w_0|\le r_0$.

Let $\xi\in K$ and $|w-w_0|\le r_0$ with $\tilde Q(\xi,w)=0$. Then $s\le v(\xi)$ by the definition of $v$. If $s=v(\xi)$, then $(\xi,(a,u))\in F\cap N$ with frozen coordinates $a$, so $u=\upsilon(\xi)$. Hence $\varphi(\xi)$ is the unique maximiser of the coordinate $s$ on the zero set of $\tilde Q(\xi,\cdot)$ in the ball, and \cref{thm:local} applies with $\lambda$ the unit vector of $s$. It gives a ball $V^*\subseteq U_1$ and a set $E^*$ of smaller dimension such that at every $\xi_1\in V^*\setminus E^*$ the function $v$ is $C^D$ and its Taylor polynomial is a limit of polynomials computed by a circuit graph that depends only on the program, on $N$, $D$ and on the set of frozen coordinates. There are finitely many such graphs.

The ball $U_1$ was arbitrary in the dense open set $A_1$, so the points $\xi_1$ obtained are dense in $A$. For each of the finitely many graphs, the set of $\xi_1\in A$ near which $v$ is $C^D$ and at which its Taylor polynomial is such a limit is defined by a first-order formula, so it is semialgebraic by \ref{S1}. Let $A^\circ$ be the union of these sets.
\end{proof}

\subsection{The permanent}

\begin{lemma}\label{lem:ryser}
The languages $\{Y=\per_m(X)\}$, $\{Y\le\per_m(X)\}$ and $\{Y\ge\per_m(X)\}$ are in $\PARR$, and $\{Y=\per_m(X)\}$ is in $\PARC$.
\end{lemma}

\begin{proof}
By Ryser's formula \cite{Rys63},
\[
\per_m(X)=\sum_{S\subseteq[m]}(-1)^{m-|S|}\prod_{i=1}^m\sum_{j\in S}X_{ij}.
\]
Use one processor for each of the $2^m$ subsets $S$ and combine the results by a binary tree of additions of depth $m$. This gives circuits with $2^m\mathrm{poly}(m)$ gates and depth $O(m)$ whose gates are described in polynomial time from their indices, followed by one comparison with $Y$.
\end{proof}

\begin{corollary}\label{cor:np}
Let trees of depth $t_m$ with $k_m$ certificate coordinates decide $\{Y=\per_m(X)\}$, $\{Y\le\per_m(X)\}$ or $\{Y\ge\per_m(X)\}$ in the existential sense of \cref{thm:graph,thm:threshold}. Then $t_m\ge2^{\Omega((\log m)^{4/3})}$. In particular none of these languages is in $\NPR$, and $\NPR\subsetneq\PARR$.
\end{corollary}

\begin{proof}
By \cref{thm:graph,thm:threshold} with $n=m^2$, $d=m$ and $\bar k\le2t_m$, the border size of $\per_m$ is $O(m^3t_m+m^2t_m^2)$. By \cref{thm:lowerbound} it exceeds $2^{c(\log m)^{4/3}}$ for large $m$. A machine witnessing membership in $\NPR$ unrolls into trees of polynomial depth. The languages are in $\PARR$ by \cref{lem:ryser}, and $\NPR\subseteq\PARR$.
\end{proof}

\section{Universal certificates and the complex field}\label{sec:universal}

If a machine places the graph of a polynomial $f$ in $\coNPR$, then every point off the graph has a rejected certificate, and no point on the graph has one. The rejected certificates of the points slightly above the graph lie in finitely many cells, and no point of the graph has a certificate in any of these cells. We measure how far such a certificate lies from the boundary of its cell and from infinity. The largest such margin is the optimal value of a polynomial program, so \cref{prop:value} applies to it. The margin tends to $0$ at the graph, and its level sets are hypersurfaces that converge to the graph together with their derivatives. The Taylor polynomials of a level set are obtained from those of the margin by a Newton iteration with one unknown. This proves \cref{thm:onesided}. Passing to real and imaginary parts gives the same conclusion over $\CC$ (\cref{thm:complex}), and with \cref{thm:lowerbound} we obtain the separations of \cref{cor:conp}.

\subsection{One-sided certificates}\label{sec:universal-onesided}

Let a tree on inputs $(x,y,z)\in\RR^n\times\RR\times\RR^k$ be given, and designate some of its leaves as \emph{marked}. A triple $(x,y,z)$ is \emph{marked} if it reaches a marked leaf. By the convention of \cref{sec:trees} every rejection happens at a leaf. In particular, when the rejecting leaves are marked, a triple that is rejected through a division by zero is marked.

\begin{theorem}[one-sided certificates]\label{thm:onesided}
Let $f\in\RR[X_1,\dots,X_n]$ have degree at most $d$ with $d\ge1$, let $U\subseteq\RR^n$ be open and nonempty, and let a tree of depth $t$ on inputs $(x,y,z)\in\RR^n\times\RR\times\RR^k$, with some of its leaves marked, be given. Suppose that for every $x\in U$
\begin{enumerate}[label=\textup{(\roman*)}]
\item no triple $(x,f(x),z)$ is marked, and
\item there is $\rho>0$, depending on $x$, such that for every $y$ with $f(x)<y<f(x)+\rho$ some triple $(x,y,z)$ is marked.
\end{enumerate}
Then the border complexity of $f$ is at most $C(d+1)^6\bigl((t+1)^2+n\bigr)$, where $C$ is an absolute constant.
\end{theorem}

With the rejecting leaves marked, the theorem applies to a tree that, for $x\in U$, accepts $(x,y,z)$ for every $z$ exactly when $y=f(x)$, and to a tree that accepts $(x,y,z)$ for every $z$ exactly when $y\le f(x)$. In both cases every certificate of a point $(x,f(x))$ is accepted, and every point $(x,y)$ with $y>f(x)$ has a rejected certificate. With the accepting leaves marked, the theorem applies to a tree that accepts $(x,y,z)$ for some $z$ exactly when $y>f(x)$. The cases with $y\ge f(x)$ in place of $y\le f(x)$ and with $y<f(x)$ in place of $y>f(x)$ follow by replacing $(y,f)$ with $(-y,-f)$.

\subsection{One cell covers a one-sided neighborhood}\label{sec:universal-cell}

In this subsection and the next three we fix $f$, $d$ and a tree of depth $t$ with marked leaves as in \cref{thm:onesided}. We may restrict $U$ to an open ball contained in it.

By \cref{lem:cells}, which holds for any set of marked leaves, the set of marked triples is a finite union of cells
\[
B=Z\cap O,\qquad Z=\{p_1=\dots=p_r=0\},\qquad O=\{q_1>0,\dots,q_m>0\},
\]
where $r,m\le2t$ and, for each cell, all the $p_i$ and $q_j$ are computed by one straight-line program with at most $8t$ operations from the input coordinates and the constants on one path. For a cell $B$ let $\Pi_B\subseteq\RR^n\times\RR$ be its projection to $(x,y)$-space. It is semialgebraic by \ref{S1}.

Let $x\in U$. By \ref{S3} each set $\{y:(x,y)\in\Pi_B\}$ is a finite union of points and intervals, so it either contains an interval $(f(x),f(x)+e)$ with $e>0$ or is disjoint from such an interval. By \textup{(ii)} the union of these finitely many sets contains $(f(x),f(x)+\rho)$. Hence at least one of them contains an initial interval $(f(x),f(x)+e)$ with $e>0$. For each cell $B$ the set of $x\in U$ for which this happens is semialgebraic by \ref{S1}, and these finitely many sets cover $U$. By \ref{S4}(a) one of them has nonempty interior. We fix the corresponding cell $B$. Certificate coordinates that do not occur in the polynomials of $B$ do not affect membership in $B$. We fix them to $0$ and let $k$ denote the number of the others, which is at most $2t$ by \cref{lem:cells}. This does not change $\Pi_B$.

For $x$ in the interior of the chosen set let $e(x)$ be the supremum of the numbers $e\in(0,1]$ with $\{x\}\times(f(x),f(x)+e)\subseteq\Pi_B$. Then $\{x\}\times(f(x),f(x)+e(x))\subseteq\Pi_B$, and $e$ is a positive semialgebraic function. By \ref{S5} it is continuous on an open dense set, so it is bounded below by some $e_0>0$ on an open ball $U_1$. Put
\[
A=\{(x,y):\ x\in U_1,\ f(x)<y<f(x)+e_0\}.
\]
This is a nonempty open semialgebraic subset of $\RR^{n+1}$ with $A\subseteq\Pi_B$. Since $B$ consists of marked triples, it follows from \textup{(i)} that no point $(x,f(x))$ with $x\in U$ lies in $\Pi_B$.

\subsection{The margin of a certificate}\label{sec:universal-margin}

A certificate in $B$ may tend to infinity, and the strict inequalities $q_j>0$ that it satisfies may approach equality, as $y$ decreases to $f(x)$. The margin measures both effects.

Substitute $z=\tilde u/u_0$, with new variables $u_0\in\RR$ and $\tilde u\in\RR^k$, in the program of $B$, and keep numerators as in \cref{sec:cells}. More precisely, we represent each value $V$ of the program by a pair $(P,R)$ of polynomials in $(x,y,u_0,\tilde u)$ such that $R=u_0^a$ for some integer $a\ge0$ and $V(x,y,\tilde u/u_0)=P/R$ whenever $u_0\ne0$. The input coordinates $x_i$ and $y$ and the constants $c$ are represented by $(x_i,1)$, $(y,1)$ and $(c,1)$, a certificate coordinate $z_i$ by $(\tilde u_i,u_0)$, and
\[
(P_1,R_1)\pm(P_2,R_2)=(P_1R_2\pm P_2R_1,\ R_1R_2),\qquad(P_1,R_1)\cdot(P_2,R_2)=(P_1P_2,\ R_1R_2).
\]
Each operation of the program is replaced by at most four operations on such pairs. We obtain polynomials $p_i^\sharp$ and $q_j^\sharp$ in $(x,y,u_0,\tilde u)$, computed together by one straight-line program with at most $32t$ operations, and integers $a_i,b_j\ge0$ such that
\[
p_i\bigl(x,y,\tilde u/u_0\bigr)=\frac{p_i^\sharp(x,y,u_0,\tilde u)}{u_0^{a_i}},\qquad q_j\bigl(x,y,\tilde u/u_0\bigr)=\frac{q_j^\sharp(x,y,u_0,\tilde u)}{u_0^{b_j}}\qquad\text{whenever }u_0\ne0 .
\]
The exponents $a_i$ and $b_j$ may be exponentially large in $t$. They are never written out, because the program computes the powers of $u_0$ as polynomials and forms sums as $P_1R_2\pm P_2R_1$.

Consider the system
\begin{equation}\label{eq:margin}
u_0^2+|\tilde u|^2=1,\qquad p_i^\sharp=0,\qquad u_0-s=\sigma_0^2,\qquad q_j^\sharp-s=\sigma_j^2\qquad(1\le i\le r,\ 1\le j\le m)
\end{equation}
in the unknowns $s\in\RR$ and $\nu=(u_0,\tilde u,\sigma_0,\dots,\sigma_m)\in\RR^\kappa$, where $\kappa=k+m+2\le4t+2$, with parameters $\xi=(x,y)\in\RR^{n+1}$. Let
\[
Q=\bigl(u_0^2+|\tilde u|^2-1\bigr)^2+\sum_{i=1}^r\bigl(p_i^\sharp\bigr)^2+\bigl(u_0-s-\sigma_0^2\bigr)^2+\sum_{j=1}^m\bigl(q_j^\sharp-s-\sigma_j^2\bigr)^2\ \in\RR[\xi,s,\nu]
\]
be the sum of the squares of its equations. Then $Q\ge0$, the zero set of $Q(\xi,\cdot)$ is the set of solutions of \eqref{eq:margin} at $\xi$, and $Q$ has a straight-line program with $\ell\le50t+8$ operations: at most $32t$ for the $p_i^\sharp$ and $q_j^\sharp$, and at most $2k+2r+5m+8\le18t+8$ for the equations, their squares and their sum.

\begin{lemma}[margin]\label{lem:margin}
\begin{enumerate}[label=\textup{(\alph*)}]
\item A point $(x,y)\in\RR^{n+1}$ lies in $\Pi_B$ exactly when \eqref{eq:margin} has a solution with $s>0$.
\item For $(x,y)\in A$ the coordinate $s$ attains a maximum $v(x,y)$ on the set of solutions of \eqref{eq:margin}, and $0<v(x,y)\le1$.
\item For every $x\in U_1$, $v(x,y)\to0$ as $y$ decreases to $f(x)$.
\end{enumerate}
\end{lemma}

We call the function $v\colon A\to(0,1]$ the \emph{margin}.

\begin{proof}
(a) Let $(s,\nu)$ be a solution with $s>0$. Then $u_0=s+\sigma_0^2\ge s>0$, $q_j^\sharp=s+\sigma_j^2\ge s>0$ and $p_i^\sharp=0$. So $z=\tilde u/u_0$ satisfies $p_i(x,y,z)=0$ and $q_j(x,y,z)>0$ for all $i$ and $j$, that is, $(x,y,z)\in B$. Conversely, let $(x,y,z)\in B$ and put $(u_0,\tilde u)=(1,z)/\sqrt{1+|z|^2}$. Then $u_0>0$ and $\tilde u/u_0=z$, so $p_i^\sharp=u_0^{a_i}p_i(x,y,z)=0$ and $q_j^\sharp=u_0^{b_j}q_j(x,y,z)>0$. With $s=\min(u_0,q_1^\sharp,\dots,q_m^\sharp)>0$, $\sigma_0=(u_0-s)^{1/2}$ and $\sigma_j=(q_j^\sharp-s)^{1/2}$ this is a solution of \eqref{eq:margin}.

(b) Let $(x,y)\in A$. The solutions of \eqref{eq:margin} with $s\ge0$ form a closed set. It is bounded: $(u_0,\tilde u)$ lies on the unit sphere, $0\le s\le s+\sigma_0^2=u_0\le1$, so that $\sigma_0^2\le1$, and $\sigma_j^2\le q_j^\sharp(x,y,u_0,\tilde u)$, which is bounded on the unit sphere. So this set is compact. Since $A\subseteq\Pi_B$, by (a) it contains a solution with $s>0$. Hence $s$ attains on it a maximum $v(x,y)>0$, which is also the maximum of $s$ on all solutions, and $v(x,y)\le u_0\le1$.

(c) Let $x\in U_1$. The function $y\mapsto v(x,y)$ on $(f(x),f(x)+e_0)$ is semialgebraic by \ref{S1} and bounded, so by \ref{S3} it has a limit $\gamma\in[0,1]$ as $y$ decreases to $f(x)$. Take $y_l$ decreasing to $f(x)$ and solutions $(v(x,y_l),\nu_l)$ of \eqref{eq:margin} at $(x,y_l)$. As in (b) they lie in a compact set, since each $q_j^\sharp(x,\cdot,\cdot)$ is bounded on the compact set $[f(x),f(x)+e_0]\times\{u_0^2+|\tilde u|^2=1\}$. A limit point is a solution of \eqref{eq:margin} at $(x,f(x))$ with $s=\gamma$. If $\gamma>0$, then $(x,f(x))\in\Pi_B$ by (a), which contradicts \textup{(i)}. So $\gamma=0$.
\end{proof}

\subsection{Level sets of the margin}\label{sec:universal-levels}

Let $x\in U_1$. By \cref{lem:margin} the function $y\mapsto v(x,y)$ on $(f(x),f(x)+e_0)$ is positive and semialgebraic, and it tends to $0$ as $y$ decreases to $f(x)$. By the monotonicity theorem \ref{S3}, applied to $h\mapsto v(x,f(x)+h)$, it is continuous and either constant or strictly monotone on some interval $(f(x),f(x)+e)$ with $e>0$. A positive constant does not tend to $0$, and neither does a positive strictly decreasing function. So $v(x,\cdot)$ is continuous and strictly increasing on such an interval. Let $e_1(x)$ be the supremum of the numbers $e\le e_0$ for which this holds. Then it holds on $(f(x),f(x)+e_1(x))$, and $e_1$ is a positive semialgebraic function on $U_1$ by \ref{S1}. By \ref{S5} there are an open ball $U_2\subseteq U_1$ and $e_2>0$ with $e_1\ge e_2$ on $U_2$. The supremum of $v(x,\cdot)$ on $(f(x),f(x)+e_2)$ is again a positive semialgebraic function of $x\in U_2$, so by \ref{S5} it is at least some $\delta_0>0$ on an open ball $U_3\subseteq U_2$.

For $x\in U_3$ and $0<\delta<\delta_0$ let $f_\delta(x)$ be the unique $y\in(f(x),f(x)+e_2)$ with $v(x,y)=\delta$. It exists by the intermediate value theorem, since $v(x,\cdot)$ is continuous and strictly increasing on this interval, with infimum $0$ and supremum at least $\delta_0$. The map $(x,\delta)\mapsto f_\delta(x)$ is semialgebraic by \ref{S1}. Since $\delta\mapsto f_\delta(x)$ is the inverse of a continuous strictly increasing function, $f_\delta(x)$ decreases to $f(x)$ as $\delta$ decreases to $0$. Thus the level sets of the margin converge pointwise to the graph of $f$.

\Cref{lem:derivatives} and \cref{cor:derivatives} require only pointwise convergence. Apply \cref{cor:derivatives} with $V=U_3$, $W(x,\delta)=f_\delta(x)$, $\varphi=f$ and $D=d$, with $x$ in the role of $\xi$ and $\delta$ in the role of $\eps$. It gives an open ball $V^*\subseteq U_3$ and $\delta_1\in(0,\delta_0)$ such that $(x,\delta)\mapsto f_\delta(x)$ is of class $C^1$ on $V^*\times(0,\delta_1)$ and has continuous $x$-derivatives of all orders up to $d$ there, and
\begin{equation}\label{eq:levelconv}
\sup_{x\in V^*}\bigl|\partial^\alpha f_\delta(x)-\partial^\alpha f(x)\bigr|\longrightarrow0\qquad(\delta\to0,\ |\alpha|\le d).
\end{equation}

\subsection{From the margin to the polynomial}\label{sec:universal-assembly}

\begin{proof}[Proof of \cref{thm:onesided}]
\emph{The margin as a value function.} Apply \cref{prop:value} to $Q$ on the open set $A$, with $\xi=(x,y)$, $N=n+1$, $\kappa=k+m+2$ and $D=d$. Its hypotheses hold: $Q$ is nonnegative and has a straight-line program with $\ell\le50t+8$ operations, and by \cref{lem:margin}(b), for every $\xi\in A$ the coordinate $s$ attains the maximum $v(\xi)$ on the set $\{(s,\nu):Q(\xi,s,\nu)=0\}$. The proposition gives a dense semialgebraic set $A^\circ\subseteq A$ and finitely many circuit graphs, each with at most
\[
L_A=C\bigl((d+1)^3(50t+8)+(d+1)^2(4t+3)^2+n+1\bigr)
\]
gates, such that for every $\xi_1\in A^\circ$ the margin $v$ is of class $C^d$ near $\xi_1$ and its Taylor polynomial of order $d$ at $\xi_1$ is a coefficientwise limit of polynomials of degree at most $d$ computed by one of these graphs with real constants.

\emph{Nondegenerate points.} Let $A'=\{(x,y):\ x\in V^*,\ f(x)<y<f(x)+e_2\}$, an open subset of $A$, and let $G$ be the set of $\xi\in A'\cap A^\circ$ with $\partial v/\partial y(\xi)\ne0$. It is semialgebraic by \ref{S1}. The set $A'\setminus A^\circ$ has empty interior, because $A^\circ$ is dense in $A$. The set of $\xi\in A'\cap A^\circ$ with $\partial v/\partial y(\xi)=0$ also has empty interior: on a ball inside it $v$ would be of class $C^1$ with $\partial v/\partial y=0$, hence constant on the segments in the ball parallel to the $y$-axis, whereas $v(x,\cdot)$ is strictly increasing on $(f(x),f(x)+e_2)$ for $x\in V^*$. By \ref{S4}(a), $\dim(A'\setminus G)\le n$.

\emph{Finitely many exceptional levels.} For $0<\delta<\delta_1$ put $H_\delta=\{(x,f_\delta(x)):\ x\in V^*\}$. Since $v(x,\cdot)$ is strictly increasing on $(f(x),f(x)+e_2)$ for $x\in V^*$, the set $H_\delta$ consists of the points of $A'$ at which $v=\delta$. So the $H_\delta$ are pairwise disjoint level sets of $v$. Let $\Delta$ be the set of $\delta\in(0,\delta_1)$ for which $H_\delta\setminus G$ has dimension $n$. It is semialgebraic by \ref{S1}, because dimension is definable in semialgebraic families. The map $\Phi(x,\delta)=(x,f_\delta(x))$ is continuous and injective on $V^*\times(0,\delta_1)$, with inverse $(x,y)\mapsto(x,v(x,y))$ on its image. Pull $A'\setminus G$ back along $\Phi$. The fiber of $\Phi^{-1}(A'\setminus G)$ over $\delta$ is $\{x\in V^*:(x,f_\delta(x))\notin G\}$, which is the image of $H_\delta\setminus G$ under the projection $(x,y)\mapsto x$, and this projection is injective on $H_\delta$. By \ref{S8} the fiber has the same dimension as $H_\delta\setminus G$. If $\Delta$ were infinite, it would contain an interval by \ref{S3}, and $\Phi^{-1}(A'\setminus G)$ would have dimension $n+1$ by \ref{S4}(c). But $\Phi$ maps this set bijectively onto a subset of $A'\setminus G$, so by \ref{S8} its dimension is at most $\dim(A'\setminus G)\le n$. Hence $\Delta$ is finite. For $\delta\in(0,\delta_1)\setminus\Delta$ the set $\{x\in V^*:(x,f_\delta(x))\notin G\}$ has dimension less than $n$, hence empty interior by \ref{S4}(a), and so the points $x\in V^*$ with $(x,f_\delta(x))\in G$ are dense in $V^*$.

\emph{A level set near the graph.} Fix $\eta>0$. Since $\Delta$ is finite, \eqref{eq:levelconv} gives $\delta\in(0,\delta_1)\setminus\Delta$ with $|\partial^\alpha f_\delta-\partial^\alpha f|\le\eta$ on $V^*$ for all $|\alpha|\le d$. Choose $x_1\in V^*$ with $\xi_1=(x_1,y_1)\in G$, where $y_1=f_\delta(x_1)$. Near $\xi_1$ the margin $v$ is of class $C^d$ with $\partial v/\partial y(\xi_1)\ne0$, and $v(x,f_\delta(x))=\delta$ for $x\in V^*$, with $f_\delta$ continuous. By the implicit function theorem the equation $v(x,y)=\delta$ has near $\xi_1$ a unique continuous solution $y$ as a function of $x$ through $\xi_1$, and this solution is of class $C^d$. So $f_\delta$ is of class $C^d$ near $x_1$. Let $P$ be the Taylor polynomial of order $d$ of $v$ at $\xi_1$, and let $g$ be that of $f_\delta$ at $x_1$. Then $g$ depends only on $P$: it is the unique polynomial of degree at most $d$ with $g(x_1)=y_1$ and
\begin{equation}\label{eq:taylorlevel}
P\bigl(x,g(x)\bigr)\equiv\delta\qquad\text{modulo terms of order above $d$ in $x-x_1$.}
\end{equation}
Indeed, by the chain rule the Taylor polynomial of order $d$ at $x_1$ of the constant function $x\mapsto v(x,f_\delta(x))$ is obtained by truncating $P(x,g(x))$ at order $d$, so $g$ satisfies \eqref{eq:taylorlevel}. If a polynomial $g'\ne g$ of degree at most $d$ with $g'(x_1)=y_1$ also satisfies it, let $h_j$ be the nonzero homogeneous part of lowest order $j$ of $g'-g$ in $x-x_1$, so that $1\le j\le d$. The part of order $j$ of $P(x,g'(x))-P(x,g(x))$ is $\partial P/\partial y(\xi_1)\,h_j$, which is nonzero because $\partial P/\partial y(\xi_1)=\partial v/\partial y(\xi_1)\ne0$. This contradicts \eqref{eq:taylorlevel} for $g$ and $g'$.

Written in powers of $x-x_1$, the coefficients of $g$ are the numbers $\partial^\alpha f_\delta(x_1)/\alpha!$ and, since $\deg f\le d$, those of $f$ are the numbers $\partial^\alpha f(x_1)/\alpha!$, for $|\alpha|\le d$. They differ by at most $\eta$. Expanding the powers of $x-x_1$ and using that $V^*$ is bounded, we obtain a constant $C_1$, depending on $n$, $d$ and $V^*$ but not on $\eta$, such that every coefficient of $g-f$ is at most $C_1\eta$ in absolute value.

\emph{The circuit.} Let $\Rtrunc=\RR[\chi][\tau]/(\tau^{d+1})$, with $\chi=(\chi_1,\dots,\chi_n)$, be the ring of \cref{lem:newton} with $D=d$, where $\chi=x-x_1$ plays the role of the increment. Put $b=\bigl(\partial P/\partial y(\xi_1)\bigr)^{-1}$, $Y^{(0)}=y_1$ and
\[
Y^{(j+1)}=Y^{(j)}-b\,\bigl(P(x_1+\tau\chi,\,Y^{(j)})-\delta\bigr)\qquad(0\le j<d).
\]
This is the iteration of \cref{lem:newton} with one unknown, for $E(x,y)=P(x,y)-\delta$, the point $x_1$ in place of $\xi_1$, $w_1=y_1$ and the $1\times1$ matrix $b$ in the role of the matrix called $A$ there. We have $E(x_1,y_1)=v(\xi_1)-\delta=0$, and $b$ is the inverse of $\partial E/\partial y(x_1,y_1)$. Let $\omega$ be the analytic function near $x_1$ with $\omega(x_1)=y_1$ and $P(x,\omega(x))=\delta$. By the chain rule its Taylor polynomial of order $d$ at $x_1$ satisfies \eqref{eq:taylorlevel}, so it equals $g$. By \cref{lem:newton}(b), $g$ is the sum of the $d+1$ coefficients of $Y^{(d)}$ with $\chi=x-x_1$ substituted; that is, $Y^{(d)}$ is the expansion of $g(x_1+\tau\chi)$.

Let $\Gamma$ be one of the circuit graphs given by \cref{prop:value}, and let $\Psi$ be a polynomial of degree at most $d$ that $\Gamma$ computes with real constants. Replace $P$ in the iteration by $\Psi$, evaluated gate by gate in $\Rtrunc$ with $x_1+\tau\chi$ and $Y^{(j)}$ fed into the variable gates of $\Gamma$, and take the sum of the coefficients of the result. This defines a composite circuit graph that depends only on $\Gamma$, $n$ and $d$. It has at most $d(d+1)(d+2)L_A+O(d^2+n)$ gates: each of the $d$ steps evaluates $\Gamma$ in $\Rtrunc$ at a cost of at most $(d+1)(d+2)$ gates per gate of $\Gamma$, the subtraction of $\delta$, the multiplication by $b$ and the subtraction from $Y^{(j)}$ cost $O(d)$ gates, and forming $\chi$ and the final sum costs $O(n+d)$ gates. Its constants are those of $\Psi$ together with $x_1$, $y_1$, $\delta$ and $b$. By \cref{lem:newton}(a), applied with $E=\Psi-\delta$, the output has degree at most $d$, and, since $\Psi$ has degree at most $d$ and evaluating a circuit in $\Rtrunc$ evaluates the polynomial it computes, the coefficients of the output are fixed polynomial functions of the coefficients of $\Psi$ and of $x_1,y_1,\delta,b$. At the coefficients of $P$ these functions give the coefficients of $g$. By the choice of $\xi_1\in A^\circ$, $P$ is a coefficientwise limit of polynomials $\Psi$ of degree at most $d$ computed by one graph $\Gamma$ with real constants. By continuity, $g$ is a coefficientwise limit of outputs of the corresponding composite graph with real constants. So one such output has all its coefficients within $C_1\eta$ of those of $g$, and hence within $2C_1\eta$ of those of $f$.

\emph{Conclusion.} For every $\eta>0$ one of finitely many composite graphs computes, with real constants, a polynomial whose coefficients are within $2C_1\eta$ of those of $f$. By \cref{lem:closure} one of these graphs approximates $f$. Since $L_A\le C'(d+1)^3\bigl((t+1)^2+n\bigr)$ and $d(d+1)(d+2)\le(d+1)^3$, its size is at most $C''(d+1)^6\bigl((t+1)^2+n\bigr)$ for absolute constants $C'$ and $C''$.
\end{proof}

\subsection{The complex field}\label{sec:universal-complex}

A \emph{complex tree} is defined as in \cref{def:tree}, with complex inputs and complex constants and with tests $v=0$ only. As over $\RR$, an input that reaches a division by zero is rejected, and every rejection happens at a leaf. A machine over $\CC$ that runs in time $T$ unrolls into a complex tree of depth $O(T)$.

\begin{theorem}[complex field]\label{thm:complex}
Let $f\in\RR[X_1,\dots,X_n]$ have degree at most $d$ with $d\ge1$. Suppose that a complex tree of depth $t$ on inputs $(x,y,z)\in\CC^n\times\CC\times\CC^k$ satisfies one of the following for all $(x,y)\in\CC^n\times\CC$:
\begin{enumerate}[label=\textup{(\alph*)}]
\item the tree accepts $(x,y,z)$ for some $z\in\CC^k$ exactly when $y=f(x)$;
\item the tree accepts $(x,y,z)$ for every $z\in\CC^k$ exactly when $y=f(x)$.
\end{enumerate}
Then the border complexity of $f$ is at most $C(d+1)^6\bigl((t+1)^2+n\bigr)$, where $C$ is an absolute constant. In particular a circuit over $\Leps\subseteq\CC((\eps))$ of this size approximates $f$.
\end{theorem}

\begin{proof}
Write every complex number as the pair of its real and imaginary parts. A complex addition, subtraction or multiplication costs a bounded number of real operations. A complex division $v_1/v_2$ costs a bounded number of real operations, among them divisions by $(\operatorname{Re}v_2)^2+(\operatorname{Im}v_2)^2$, which vanishes exactly when $v_2=0$. A test $v=0$ becomes the two real tests $\operatorname{Re}v=0$ and $\operatorname{Im}v=0$, and a complex constant becomes two real constants. This turns the complex tree into a real tree of depth at most $c_0t$, for an absolute constant $c_0$, that accepts a real input exactly when the complex tree accepts the corresponding complex input; its $2k$ certificate coordinates are the real and imaginary parts of $z$. Replace the inputs $\operatorname{Im}x$ and $\operatorname{Im}y$ by the constant $0$. This gives a real tree $T'$ of depth at most $c_0t$ on inputs $(x,y,z')\in\RR^n\times\RR\times\RR^{2k}$. Since $f$ has real coefficients, for real $x$ and $y$ the complex equation $y=f(x)$ is the real one.

In case (a), for all $x\in\RR^n$ and $y\in\RR$ the tree $T'$ accepts $(x,y,z')$ for some $z'$ exactly when $y=f(x)$. So \cref{thm:graph} applies with $U=\RR^n$ and $k'\le2c_0t$, and it shows that the border complexity of $f$ is at most $C\bigl((d+1)^3c_0t+(d+1)^2(2c_0t+2)^2+n\bigr)$. In case (b), mark the rejecting leaves of $T'$. For real $x$ no triple $(x,f(x),z')$ is marked, because every certificate of $(x,f(x))$ is accepted, and for every real $y\ne f(x)$ some triple $(x,y,z')$ is marked. So \cref{thm:onesided} applies with $U=\RR^n$, and it shows that the border complexity of $f$ is at most $C(d+1)^6\bigl((c_0t+1)^2+n\bigr)$. Both bounds are at most a constant multiple of $(d+1)^6\bigl((t+1)^2+n\bigr)$.
\end{proof}

\begin{remark}[complex coefficients]\label{rem:complexcoef}
Let $f$ have complex coefficients, and write $f=f_1+if_2$ with $f_1,f_2\in\RR[X_1,\dots,X_n]$. In case (a) of \cref{thm:complex} the argument still applies if both the real and the imaginary part of $y$ are kept as inputs: the real tree then decides, in the existential sense, the graph of the polynomial map $(f_1,f_2)\colon\RR^n\to\RR^2$, and \cref{rem:several} with $\mu=2$ bounds the border complexity of $f_1$ and of $f_2$. In case (b) the argument does not apply: the graph of $f$ over $\RR^n$ then has real codimension two, and \cref{thm:onesided} requires a hypersurface with two sides.
\end{remark}

\subsection{The separations}\label{sec:universal-separations}

A tree \emph{decides} a language $L$ on inputs $(X,Y)$ \emph{in the universal sense} if $(X,Y)\in L$ exactly when the tree accepts $(X,Y,z)$ for every $z$, and \emph{in the existential sense} if $(X,Y)\in L$ exactly when it accepts $(X,Y,z)$ for some $z$, as in \cref{thm:graph}.

\begin{corollary}\label{cor:conp}
Suppose that for every $m$ a tree of depth $t_m$ decides one of the languages $\{Y=\per_m(X)\}$, $\{Y\le\per_m(X)\}$ and $\{Y\ge\per_m(X)\}$ in the universal sense, or one of the languages $\{Y<\per_m(X)\}$ and $\{Y>\per_m(X)\}$ in the existential sense, or that a complex tree of depth $t_m$ decides $\{Y=\per_m(X)\}\subseteq\CC^{m^2+1}$ in the existential or in the universal sense. Then $t_m\ge2^{\Omega((\log m)^{4/3})}$. In particular:
\begin{enumerate}[label=\textup{(\alph*)}]
\item the graph $\{Y=\per_m(X)\}$ and the thresholds $\{Y\le\per_m(X)\}$ and $\{Y\ge\per_m(X)\}$ of the real permanent are in neither $\NPR$ nor $\coNPR$, and the languages $\{Y<\per_m(X)\}$ and $\{Y>\per_m(X)\}$ are not in $\NPR$;
\item the graph of the complex permanent is in neither $\NPC$ nor $\coNPC$;
\item $\NPR\cup\coNPR\subsetneq\PARR$ and $\NPC\cup\coNPC\subsetneq\PARC$.
\end{enumerate}
\end{corollary}

\begin{proof}
Let $n=m^2$, $d=m$ and $U=\RR^{m^2}$. For a real tree in the universal sense, mark its rejecting leaves; for a real tree in the existential sense, mark its accepting leaves. As explained after \cref{thm:onesided}, the hypotheses of that theorem then hold with $f=\per_m$ for $\{Y=\per_m(X)\}$, $\{Y\le\per_m(X)\}$ and $\{Y>\per_m(X)\}$, and, after $Y$ is replaced by $-Y$, which adds one node to the depth, with $f=-\per_m$ for $\{Y\ge\per_m(X)\}$ and $\{Y<\per_m(X)\}$. A circuit that approximates $-\per_m$ gives one that approximates $\per_m$ with two more gates. For a complex tree, \cref{thm:complex} applies with $f=\per_m$, since the permanent has real coefficients. In every case the border complexity of $\per_m$ is $O\bigl(m^6(t_m+m)^2\bigr)$. By \cref{thm:lowerbound} it exceeds $2^{c(\log m)^{4/3}}$ for $m\ge m_0$. Hence $t_m\ge2^{\Omega((\log m)^{4/3})}$.

(a) A polynomial-time machine unrolls into trees of polynomial depth. If the graph or one of the two thresholds were in $\coNPR$, its complement would be in $\NPR$, and a machine for the complement would unroll into trees of polynomial depth that accept $(X,Y,z)$ for some $z$ exactly when $(X,Y)$ is not in the language. Exchanging their accepting and rejecting leaves, which is possible because every rejection happens at a leaf, gives trees of the same depth that decide the language in the universal sense. This contradicts the bound. A machine that places $\{Y<\per_m(X)\}$ or $\{Y>\per_m(X)\}$ in $\NPR$ unrolls into trees of polynomial depth that decide it in the existential sense, which again contradicts the bound; these two languages are the complements of the thresholds $\{Y\ge\per_m(X)\}$ and $\{Y\le\per_m(X)\}$. The graph and the thresholds are not in $\NPR$ by \cref{cor:np}, which follows from \cref{thm:graph,thm:threshold}.

(b) A machine that places the graph of the complex permanent in $\NPC$ unrolls into complex trees of polynomial depth that decide it in the existential sense. If the graph were in $\coNPC$, exchanging the accepting and rejecting leaves of the complex trees for its complement would give complex trees of polynomial depth that decide it in the universal sense. Both contradict the bound.

(c) By \cref{lem:ryser} the graph and the thresholds of the real permanent are in $\PARR$, and the graph of the complex permanent is in $\PARC$. We have $\NPR\subseteq\PARR$ by \cref{sec:trees}, and the inclusions $\coNPR\subseteq\PARR$ and $\NPC\subseteq\PARC$ are standard \cite{BCSS98}. Since $\PARC$ is closed under complements, also $\coNPC\subseteq\PARC$. So $\NPR\cup\coNPR\subseteq\PARR$ and $\NPC\cup\coNPC\subseteq\PARC$, and by (a) and (b) the graph of the permanent shows that both inclusions are strict.
\end{proof}

\section{An existential and a universal block}\label{sec:second}

\begin{theorem}\label{thm:second}
Let $f\in\RR[X_1,\dots,X_n]$ have degree at most $d$ with $d\ge1$, and let $U\subseteq\RR^n$ be open and nonempty. Suppose that a tree of depth $t$ on inputs $(x,y,z,u)\in\RR^n\times\RR\times\RR^k\times\RR^l$ satisfies, for all $x\in U$ and $y\in\RR$,
\[
\text{there is }z\text{ such that the tree accepts }(x,y,z,u)\text{ for all }u\quad\Longleftrightarrow\quad y=f(x).
\]
Then the border complexity of $f$ is at most a polynomial in $t$, $k$, $l$, $d$ and $n$.
\end{theorem}

Once a certificate $z$ for which every $u$ is accepted has been isolated, the points $(x,y,z)$ at which every $u$ is accepted form a graph $\Gamma$ of lower dimension, and every other nearby point has a rejecting $u$. Rejection is the existential condition, so no tested polynomial need vanish on $\Gamma$. The largest margin $v$ of a rejecting certificate is positive off $\Gamma$ and zero on it. However, $v$ is defined as a maximum. It need not be differentiable at the points of $\Gamma$, and it can be discontinuous on sets of lower dimension, so it can take small values far from $\Gamma$. We therefore replace $v$ by an envelope that is the value of a single polynomial program, is continuously differentiable, and vanishes exactly on $\Gamma$.

Throughout the section the hypotheses of \cref{thm:second} are in force.

\subsection{Isolating a certificate}

Let $\mathcal T$ be the set of $(x,y,z)$ with $x\in U$ at which every $u$ is accepted. It is semialgebraic by \ref{S1}, and by hypothesis its projection to $(x,y)$ is the graph of $f$ over $U$. Put $F=\{(x,z)\in U\times\RR^k:(x,f(x),z)\in\mathcal T\}$, which projects onto $U$. Apply \cref{lem:isolate} to $F$ over a ball in $U$ and fix coordinates. We write $z=(z^\flat,z')$, where $z^\flat$ is the vector of fixed values and $z'\in\RR^{k'}$, and $\zeta(x)=(z^\flat,\zeta'(x))$. Put
\[
w=(y,z')\in\RR^\kappa,\qquad\kappa=1+k',\qquad\varphi(x)=\bigl(f(x),\zeta'(x)\bigr),\qquad w_0=\varphi(x_0).
\]

\begin{lemma}\label{lem:goodz}
There are $r_0>0$ and a closed ball $K$ about $x_0$, inside the ball $U_3$ of \cref{lem:isolate}, such that $|\varphi(x)-w_0|\le r_0/2$ on $K$ and, for $x\in K$ and $w$ in the closed ball $\bar B$ of radius $r_0$ about $w_0$,
\[
\text{some }u\text{ is rejected at }(x,y,(z^\flat,z'))\quad\Longleftrightarrow\quad w\ne\varphi(x).
\]
\end{lemma}

\begin{proof}
Choose $K$ and $r_0$ such that the bound on $\varphi$ holds and $(x,(z^\flat,z'))$ lies in the neighborhood of $q_0$ given by \cref{lem:isolate} whenever $x\in K$ and $w\in\bar B$. If no $u$ is rejected at $(x,y,(z^\flat,z'))$, this point lies in $\mathcal T$, and $y=f(x)$ by the hypothesis. Hence $(x,(z^\flat,z'))$ is a point of $F$ in that neighborhood with fixed coordinates $z^\flat$, and therefore $z'=\zeta'(x)$. Conversely, the point $(x,f(x),\zeta(x))$ lies in $\mathcal T$.
\end{proof}

Let $\Gamma=\{(x,\varphi(x)):x\in K\}$. From now on the fixed coordinates are constants, and we regard the tree as a tree on inputs $(x,w,u)$.

\subsection{The margin and a power of it}

Apply \cref{lem:cells} with the rejecting leaves marked. For each cell consider the margin system \eqref{eq:margin} of \cref{sec:universal}, with parameters $(x,w)$ in place of $(x,y)$ and unknowns $(\nu,s)$. The proof of \cref{lem:margin}(a) applies to any cell and any parameters, so this system has a solution with $s>0$ exactly when $(x,w)$ has a rejecting certificate in that cell. The system of a cell involves only the coordinates of $u$ that occur in the cell, and there are at most $2t$ of them. Embed the unknowns $\nu$ of the different cells in a common space, with the additional coordinates set to $0$. Let $\mathcal Y$ be the set of all $(x,w,\nu,s)$ with $x\in K$, $w\in\bar B$ and $s\ge0$ that solve the system of some cell. It is closed, and its fibers over $(x,w)$ are compact and uniformly bounded: $\nu$ consists of a point of the unit sphere and of slack variables whose squares are bounded on $K\times\bar B$, and $0\le s\le1$. Define
\[
v(x,w)=\max\{s:(x,w,\nu,s)\in\mathcal Y\}.
\]

\begin{lemma}\label{lem:flatten}
\begin{enumerate}[label=\textup{(\alph*)}]
\item The fiber of $\mathcal Y$ over every $(x,w)\in K\times\bar B$ is nonempty, $0\le v\le1$, and $v(x,w)=0$ exactly when $w=\varphi(x)$.
\item $v$ is upper semicontinuous.
\item There are an integer $N\ge1$ and a real $a>0$ such that
\[
v(x,w)^N\le a\,|w-\varphi(x)|^2\qquad\text{for all }x\in K,\ w\in\bar B .
\]
\item The integer $N$ can be chosen with $\log N$ bounded by a polynomial in $t$, $n$, $k$ and $l$.
\end{enumerate}
\end{lemma}

\begin{proof}
(b) Let $(x_j,w_j)\to(x,w)$ with $v(x_j,w_j)\to\lambda$. Take solutions with $s=v(x_j,w_j)$ and pass to a convergent subsequence of solutions of the system of a single cell. The limit solves that cell's system at $(x,w)$ with $s=\lambda$, so $v(x,w)\ge\lambda$.

(a) Off $\Gamma$ a rejecting certificate exists by \cref{lem:goodz}, so $v>0$. On $\Gamma$ there is no rejecting certificate, hence no solution with $s>0$. Since $\varphi(x)$ is an interior point of $\bar B$, there are $w_j\to\varphi(x)$ with $w_j\ne\varphi(x)$, and a limit of solutions at the points $(x,w_j)$ is a solution at $(x,\varphi(x))$, necessarily with $s=0$. So the fiber is nonempty and $v=0$ there.

(c) Let $g(\sigma)$ be the supremum of $v(x,w)$ over $x\in K$ and $w\in\bar B$ with $|w-\varphi(x)|\le\sigma$. It is semialgebraic and nondecreasing. If $g(\sigma)\ge\lambda>0$ for all $\sigma>0$, a limit point of approximate maximizers would lie on $\Gamma$ and have $v\ge\lambda$ by (b), contrary to (a). So $g(\sigma)\to0$ as $\sigma\to0$, and by \ref{S7} there are $L$ and a rational $\gamma>0$, the leading exponent of $g$, with $g(\sigma)\le L\sigma^\gamma$ for small $\sigma$. Since $v\le1$, the bound $v(x,w)\le L'|w-\varphi(x)|^\gamma$ holds on all of $K\times\bar B$ for a larger $L'$. Take $N\ge2/\gamma$ and $a=L'^N(2r_0)^{N\gamma-2}$.

(d) The function $g$ is defined by a first-order formula with a bounded number of quantifier blocks over the tests of the tree: at most $2^{O(t)}$ polynomials of degree $2^{O(t)}$ in $O(n+k+l+t)$ variables. By effective quantifier elimination \cite[Chapter~14]{BPR06} its graph is defined by a quantifier-free formula in which the polynomials are at most $2^{q}$ in number and of degree at most $2^{q}$, with $q$ polynomial in $t,n,k,l$. The graph has empty interior in $\RR^2$, so each of its points is a zero of one of these polynomials, and all but finitely many of its points are zeros of one that has positive degree in the second variable. Let $P$ be the product of the latter. Then $P(\sigma,g(\sigma))=0$ for small $\sigma>0$, and the leading exponent $\gamma$ of the branch $g$ is a slope of the Newton polygon of $P$, a ratio of positive integers not exceeding $\deg P\le2^{2q}$. So $\gamma\ge2^{-2q}$, and $N=\lceil2/\gamma\rceil$ has $\log N\le2q+2$.
\end{proof}

Put $\hat v=v^N$. Below we use only the upper semicontinuity of $v$ and its zero set. The function $v$ may be discontinuous, and its small values may occur far from $\Gamma$.

\subsection{The envelope}

Fix $x\in K$. An \emph{atom} is a point $(\omega,\nu,s)$ with $(x,\omega,\nu,s)\in\mathcal Y$. Let $r=\kappa+3$. For $c\in\RR^\kappa$ let $e(x;c)$ be the maximum, over weights $\lambda_1,\dots,\lambda_r\ge0$ with $\sum_i\lambda_i=1$ and atoms $(\omega_i,\nu_i,s_i)$, of
\begin{equation}\label{eq:J}
J=\sum_i\lambda_is_i^N-a\sum_i\lambda_i|\omega_i-\bar\omega|^2-2a\,|c-\bar\omega|^2,\qquad\bar\omega=\sum_i\lambda_i\omega_i .
\end{equation}
The maximum exists because the weights and atoms range over a compact set. In terms of the moment vector $m=\sum_i\lambda_i\bigl(s_i^N,|\omega_i|^2,\omega_i\bigr)=(m_0,m_2,m_1)\in\RR^{\kappa+2}$,
\[
J=j(c;m)=m_0-a\,m_2-a\,|m_1|^2+4a\langle c,m_1\rangle-2a\,|c|^2 ,
\]
because $\sum_i\lambda_i|\omega_i-\bar\omega|^2=m_2-|m_1|^2$. By Carath\'eodory's theorem the moment vectors obtained with $r$ atoms are exactly the points of the convex hull $M_x$ of the compact set of vectors $(s^N,|\omega|^2,\omega)$ of single atoms. So $e(x;c)=\max_{m\in M_x}j(c;m)$, where $M_x$ is compact and convex and $j$ is concave in $m$ and strictly concave in $m_1$. In terms of measures, $e(x;c)$ is the largest value of $\int\hat v\,d\mu-a\Var(\mu)-2a|c-\operatorname{mean}(\mu)|^2$ over probability measures $\mu$ on $\bar B$. It is a regularization of $\hat v$ of the kind studied by Lasry and Lions \cite{LL86}.

\begin{lemma}\label{lem:envelope}
For every $x\in K$:
\begin{enumerate}[label=\textup{(\alph*)}]
\item $e(x;\cdot)$ is continuously differentiable on $\RR^\kappa$, with $\nabla_ce(x;c)=4a\,(m_1^*-c)$, where $m_1^*$ is the component $m_1$ common to all maximizers. On $\bar B$ it is Lipschitz with constant $8ar_0$.
\item $e(x;c)\ge v(x,c)^N$ for $c\in\bar B$.
\item $e(x;\varphi(x))=0$.
\end{enumerate}
Hence on $\bar B$ the function $e(x;\cdot)$ is nonnegative and vanishes exactly at $\varphi(x)$.
\end{lemma}

\begin{proof}
(a) The function $e(x;c)+2a|c|^2=\max_{m\in M_x}\bigl[m_0-am_2-a|m_1|^2+4a\langle c,m_1\rangle\bigr]$ is a maximum of affine functions of $c$ over a compact set, so it is convex. By strict concavity, two maximizers with different $m_1$ would have a midpoint in $M_x$ with a larger value. So all maximizers have the same component $m_1^*$. The subdifferential of a maximum of affine functions over a compact set is the convex hull of the gradients $4am_1$ at the maximizers, which here is a single point. A convex function with a single subgradient at every point is continuously differentiable. Since $m_1^*\in\bar B$, the gradient has norm at most $8ar_0$ on $\bar B$.

(b) Take a single atom at $\omega=c$ with $s=v(x,c)$. Then $\bar\omega=c$ and $J=v(x,c)^N$.

(c) Write $\varphi=\varphi(x)$. By \cref{lem:flatten}(c), $s_i^N\le v(x,\omega_i)^N\le a|\omega_i-\varphi|^2$, and $\sum_i\lambda_i|\omega_i-\varphi|^2=\sum_i\lambda_i|\omega_i-\bar\omega|^2+|\bar\omega-\varphi|^2$. So at $c=\varphi$
\[
J\le a\,|\bar\omega-\varphi|^2-2a\,|\varphi-\bar\omega|^2\le0 ,
\]
and $e(x;\varphi)\ge0$ by (b). The last assertion follows from (b), (c) and \cref{lem:flatten}(a).
\end{proof}

Part (c) uses the power $N$. Without it, $e(x;\varphi(x))$ would be positive, and the minimizers of $e$ could move away from the graph toward the points where $v$ is discontinuous.

\begin{lemma}\label{lem:continuity}
There is a closed ball $K'\subseteq K$ such that $e$ is continuous on $K'\times\bar B$. Consequently, for every $\eta>0$,
\[
\mu(\eta)=\min\{e(x;c):x\in K',\ c\in\bar B,\ |c-\varphi(x)|\ge\eta\}>0 .
\]
\end{lemma}

\begin{proof}
The function $e$ is semialgebraic by \ref{S1}. By \ref{S5}, applied on the open set $\interior K\times\interior\bar B$, the set of its points of discontinuity there has dimension at most $n+\kappa-1$. Suppose $e$ is discontinuous at $(x^\circ,c^\circ)$ with $x^\circ\in\interior K$ and $c^\circ\in\bar B$: there are $(x_j,c_j)\to(x^\circ,c^\circ)$ with $|e(x_j;c_j)-e(x^\circ;c^\circ)|\ge\eta$. By the Lipschitz bound of \cref{lem:envelope}(a), which is uniform in $x$, we get $|e(x_j;c)-e(x^\circ;c)|\ge\eta/4$ for all large $j$ and all $c\in\bar B$ within $\eta/(32ar_0)$ of $c^\circ$. So $e$ is discontinuous at $(x^\circ,c)$ for all those $c$, which include an open subset of $\interior\bar B$. If the set of such $x^\circ$ had dimension $n$, the set of points of discontinuity would have dimension $n+\kappa$ by \ref{S4}(c). So it has smaller dimension, and $K'$ can be chosen disjoint from its closure. The minimum is then positive by compactness and \cref{lem:envelope}.
\end{proof}

\subsection{The envelope is a value function}

\begin{lemma}\label{lem:envelopevalue}
Let $\mathcal O=\interior K'\times\interior\bar B$. For every $D\ge1$ there are finitely many circuit graphs, each with at most
\[
s_e(D)=C\bigl((D+1)^3\kappa\,(t+\kappa+\log N)+(D+1)^2\kappa^2(t+\kappa)^2+n\bigr)
\]
gates, and a dense semialgebraic set $\mathcal O^\circ\subseteq\mathcal O$, such that at every $\xi_1\in\mathcal O^\circ$ the function $e$ is of class $C^D$ near $\xi_1$ and its Taylor polynomial of order $D$ at $\xi_1$ is a coefficientwise limit of polynomials of degree at most $D$ computed by one of these graphs with real constants.
\end{lemma}

\begin{proof}
Assign a rejecting cell to each of the $r$ atoms. For one assignment the constraints on the weights and atoms form one polynomial system with parameters $\xi=(x,c)$: the margin system \eqref{eq:margin} of the assigned cell for each atom, with $(x,\omega_i)$ in place of $(x,y)$; the equation $s_i=\theta_i^2$, which enforces $s_i\ge0$; the equation $|\omega_i-w_0|^2+\rho_i^2=r_0^2$, which keeps $\omega_i$ in $\bar B$; the equations $\lambda_i=\beta_i^2$ and $\sum_i\beta_i^2=1$; and $T=J$ for a new variable $T$, with $J$ as in \eqref{eq:J}. Let $Q$ be the sum of the squares of these equations. It has a straight-line program with $O(\kappa(t+\kappa+\log N))$ operations, where $s_i^N$ is computed by repeated squaring, and $O(\kappa(t+\kappa))$ unknowns. Its zero set over each $\xi$ is compact, since $0\le s_i\le1$ and the slack variables are bounded. The function $e$ is the largest, over the finitely many assignments, of the maxima of $T$ on these zero sets.

The sets on which a given assignment attains $e$ are semialgebraic and cover $\mathcal O$, so the union of their interiors is dense by \ref{S4}(a). On each interior the maximum of $T$ for that assignment is attained and equals $e$, and \cref{prop:value} applies with the coordinate $T$ in the role of $s$. The union of the dense sets and the finitely many graphs so obtained have the stated properties.
\end{proof}

Let $\mathcal P_D$ be the closure, coefficientwise, of the set of polynomials of degree at most $D$ computed by the graphs of \cref{lem:envelopevalue} with real constants.

\subsection{A generic linear term}

By \ref{S8} there is an open dense semialgebraic set $\mathrm{Reg}\subseteq\mathcal O$ on which $e$ is Nash in $(x,c)$; put $\mathrm{Sing}=\mathcal O\setminus\mathrm{Reg}$, of dimension at most $n+\kappa-1$. Consider the semialgebraic map $G(x,c)=(x,\nabla_ce(x;c))$ on $\mathcal O$, which is defined at every point by \cref{lem:envelope}(a). Let $\mathrm{Bad}\subseteq\RR^n\times\RR^\kappa$ be the union of $G(\mathrm{Sing})$ and of the set of critical values of $G$ on $\mathrm{Reg}$. By \ref{S8} and \ref{S6}, $\dim\mathrm{Bad}\le n+\kappa-1$.

\begin{lemma}\label{lem:tilt}
There are $\hat a\in\RR^\kappa$, a ball $V\subseteq\interior K'$ and $\vartheta_1>0$ such that for $x\in V$ and $0<\vartheta<\vartheta_1$ every $c\in\interior\bar B$ with $\nabla_ce(x;c)=\vartheta\hat a$ has $(x,c)\in\mathrm{Reg}$ and $\nabla_c^2e(x;c)$ invertible.
\end{lemma}

\begin{proof}
Under the semialgebraic bijection $(x,\vartheta,\hat a)\mapsto(x,\vartheta,\vartheta\hat a)$, the set of $(x,\vartheta,\hat a)$ with $\vartheta>0$ and $(x,\vartheta\hat a)\in\mathrm{Bad}$ corresponds to $\{(x,\vartheta,b):\vartheta>0,\ (x,b)\in\mathrm{Bad}\}$, which has dimension at most $n+\kappa$. By \ref{S4}(c) some $\hat a$ has a fiber of dimension at most $n$ in $(x,\vartheta)$-space. As in the first step of the proof of \cref{lem:derivatives}, there is a box $V\times(0,\vartheta_1)$ that is disjoint from this fiber. For $(x,\vartheta)$ in the box, $(x,\vartheta\hat a)\notin\mathrm{Bad}$: no point of $\mathrm{Sing}$ over $x$ has gradient $\vartheta\hat a$, and $\vartheta\hat a$ is a regular value of $\nabla_ce(x;\cdot)$ on $\mathrm{Reg}$.
\end{proof}

Fix $\hat a$, $V$ and $\vartheta_1$ as in the lemma. For $x\in K'$ and $\vartheta>0$ let $\bar c$ minimize $e(x;c)-\vartheta\langle\hat a,c\rangle$ over $c\in\bar B$. Comparison with the value at $c=\varphi(x)$, where $e=0$, gives
\[
e(x;\bar c)\le\vartheta\langle\hat a,\bar c-\varphi(x)\rangle\le2r_0|\hat a|\,\vartheta .
\]
By \cref{lem:continuity}, $|\bar c-\varphi(x)|<\eta$ as soon as $2r_0|\hat a|\vartheta<\mu(\eta)$. So the minimizers converge to $\varphi$ uniformly on $K'$ as $\vartheta\to0$. Decrease $\vartheta_1$ so that for $\vartheta<\vartheta_1$ they are interior points of the ball. Then $\nabla_ce(x;\bar c)=\vartheta\hat a$.

By \ref{S2} choose a semialgebraic map $\Theta\colon V\times(0,\vartheta_1)\to\bar B$ such that each value $\Theta(x,\vartheta)$ is a minimizer of this kind. By \cref{cor:derivatives} there are a ball $V^*\subseteq V$ and $\vartheta^*\in(0,\vartheta_1)$ such that $\Theta$ is $C^1$ on $V^*\times(0,\vartheta^*)$ and the $x$-derivatives up to order $d$ of its first coordinate converge to those of $f$, uniformly on $V^*$.

\subsection{Proof of \texorpdfstring{\cref{thm:second}}{the theorem}}

Fix $\eta>0$. Choose $\vartheta<\vartheta^*$ such that the derivatives up to order $d$ of the first coordinate of $\Theta(\cdot,\vartheta)$ are within $\eta$ of those of $f$ on $V^*$. Take $x_1\in V^*$ and put $c_1=\Theta(x_1,\vartheta)$ and $\xi_1=(x_1,c_1)$. By \cref{lem:tilt}, $\xi_1\in\mathrm{Reg}$ and $H=\nabla_c^2e(\xi_1)$ is invertible. By the implicit function theorem there is an analytic map $\psi$ near $x_1$ with $\psi(x_1)=c_1$ and $\nabla_ce(x;\psi(x))=\vartheta\hat a$, and $\Theta(\cdot,\vartheta)$, a continuous solution through $\xi_1$, agrees with it near $x_1$. So the Taylor polynomial $h$ of order $d$ at $x_1$ of the first coordinate of $\psi$ has all coefficients within $C_1\eta$ of those of $f$, where $C_1$ depends only on $n$, $d$ and $V^*$.

Let $P$ be the Taylor polynomial of $e$ of order $d+1$ at $\xi_1$. Since $e$ is Nash near $\xi_1$, its Taylor polynomials of order $d+1$ depend continuously on the base point there, and $\mathcal O^\circ$ is dense, so $P\in\mathcal P_{d+1}$ by \cref{lem:envelopevalue}. The Taylor polynomial of order $d$ of $\psi$ at $x_1$ depends only on $P$. Indeed $\nabla_cP(\xi_1)=\vartheta\hat a$ and $\nabla_c^2P(\xi_1)=H$, so there is an analytic map $\omega$ near $x_1$ with $\omega(x_1)=c_1$ and $\nabla_cP(x,\omega(x))=\vartheta\hat a$. Since $e-P$ vanishes to order $d+2$ at $\xi_1$, the difference $\nabla_ce-\nabla_cP$ vanishes to order $d+1$, and so does $\nabla_cP(x,\psi(x))-\vartheta\hat a$ at $x_1$. If $\psi-\omega$ had a nonzero homogeneous part of lowest degree $j\le d$ at $x_1$, then $\nabla_cP(x,\psi(x))-\nabla_cP(x,\omega(x))$ would have as its lowest part the image of that part under $H$, of degree $j$, which is impossible. So $\psi$ and $\omega$ have the same Taylor polynomial of order $d$ at $x_1$. That of $\omega$ is obtained from \cref{lem:newton}, with $\xi=x$, $w=c$, $D=d$,
\[
E(x,c)=\nabla_cP(x,c)-\vartheta\hat a
\]
and $A=H^{-1}$. Replace $P$ by the output $\Psi$ of one of the circuit graphs of \cref{lem:envelopevalue} for $D=d+1$, with its gradient in $c$ obtained from the circuit by \cref{thm:baurstrassen}. By \cref{lem:newton}(a) this defines a composite circuit graph with $O\bigl(d^3(s_e(d+1)+\kappa^2)+n\bigr)$ gates, whose output has degree at most $d$ and is a polynomial function of the coefficients of $\Psi$ and of $x_1$, $c_1$, $\vartheta\hat a$ and $H^{-1}$. By continuity, $h$ is a coefficientwise limit of outputs of one of the finitely many composite graphs.

So for every $\eta>0$ one of finitely many circuit graphs computes, with real constants, a polynomial of degree at most $d$ within $2C_1\eta$ of $f$. By \cref{lem:closure} one of them approximates $f$. Its size is polynomial in $t$, $k$, $d$, $n$ and $\log N$, since $\kappa\le k+1$, and $\log N$ is polynomial in $t$, $n$, $k$ and $l$ by \cref{lem:flatten}(d). \qed

\begin{corollary}\label{cor:second}
Let trees of depth $t_m$, with $k_m$ and $l_m$ coordinates in $z$ and $u$, decide $Y=\per_m(X)$ in the sense of \cref{thm:second}. Then $t_m+k_m+l_m\ge2^{\Omega((\log m)^{4/3})}$, where the constant in the exponent depends on the degree of the polynomial bound of \cref{thm:second}. In particular the graph of the permanent is not in $\Sigma_2^{\RR}$, and $\Sigma_2^{\RR}\subsetneq\PARR$.
\end{corollary}

\begin{proof}
By \cref{thm:second} with $n=m^2$ and $d=m$ the border complexity of $\per_m$ is at most a polynomial in $m$, $t_m$, $k_m$ and $l_m$. By \cref{thm:lowerbound} it exceeds $2^{c(\log m)^{4/3}}$. A $\Sigma_2^{\RR}$ machine unrolls into trees with polynomial $t_m$, $k_m$ and $l_m$. The graph is in $\PARR$ by \cref{lem:ryser}, and $\Sigma_2^{\RR}\subseteq\PARR$.
\end{proof}

\section{Interaction}\label{sec:interaction}

This section concerns interactive proofs in which the verifier is a real machine and the randomness is discrete. When the prover's messages are Boolean, a verifier for the graph of the permanent needs time $2^{\Omega((\log m)^{4/3})}$, whatever the number of rounds, and this separates the resulting class from $\PARR$. When the prover's messages are real and the coins are public, a constant number of rounds gives exactly $\NPR$, and a polynomial-time verifier for the graph of the permanent needs $\Omega(\log\log m)$ rounds.

\subsection{Boolean messages}

\begin{definition}\label{def:bip}
A \emph{verifier tree} on inputs $w\in\RR^N$ is a tree as in \cref{def:tree} with three further kinds of nodes, none of which computes a value. A \emph{coin node} has two children. A \emph{send node} is labelled by a bit, which the verifier sends to the prover, and has one child. A \emph{receive node} has two children, and the verifier continues with the child given by a bit that it receives from the prover. The \emph{transcript} of an initial segment of a path is the sequence of bits sent and received on it, in order, each marked with its direction. A \emph{prover} is a map $\mathcal P$ that assigns a bit $\mathcal P(w,s)$ to every input $w$ and every transcript $s$. Let $t$ be the depth of the tree and let $r\in\{0,1\}^t$ be a coin tape. Given $w$, $\mathcal P$ and $r$, the verifier follows the path from the root that proceeds as in \cref{def:tree} at computation and test nodes, continues with the branch $r_j$ at the $j$-th coin node on the path, and continues with the branch $\mathcal P(w,s)$ at a receive node, where $s$ is the transcript of the path up to that node. The \emph{acceptance probability} of $\mathcal P$ at $w$ is the probability that this path ends at an accepting leaf when $r$ is uniform in $\{0,1\}^t$. A \emph{Boolean-message protocol}, given by a verifier tree, \emph{decides} a language $L$ on a set $W\subseteq\RR^N$ if for every $w\in W\cap L$ some prover has acceptance probability at least $2/3$ at $w$, and for every $w\in W\setminus L$ every prover has acceptance probability at most $1/3$ at $w$.
\end{definition}

The prover is computationally unbounded and may depend on the input in an arbitrary way. It sees the bits that the verifier sends but not the coin tape, so the coins are private. A randomised prover is a probability distribution over provers, and its acceptance probability is the corresponding average of theirs; so allowing randomised provers changes neither condition in \cref{def:bip}.

A \emph{verifier} is a machine with additional instructions to flip a coin, to send a bit and to receive a bit. If it runs in time $t$ on inputs of size $N$, for every coin sequence and every sequence of received bits, then on these inputs it unrolls into a verifier tree of depth $O(t)$ whose constants are the constants of the machine. When the machine sends a computed value $v\in\{0,1\}$, the tree tests $v=0$ and continues with a send node labelled $0$ or $1$; this at most doubles the depth. The class $\BIPR$ of Ivanov and de~Rougemont \cite{IdR99} consists of the languages $L\subseteq\RR^\infty$ for which some verifier running in polynomial time decides $L$ on all of $\RR^N$, for every $N$. The number of rounds is bounded only by the running time. Ivanov and de~Rougemont proved
\[
\BIPR\subseteq\PARR
\]
in the Blum--Shub--Smale model with multiplication that is used throughout this paper \cite{IdR99}.

\begin{lemma}[boundary test]\label{lem:boolean}
Let $f\in\RR[X_1,\dots,X_n]$, let $U\subseteq\RR^n$ be open and nonempty, and let $W\subseteq\RR^n\times\RR$ be an open set that contains $(x,f(x))$ for every $x\in U$. Suppose that a Boolean-message protocol whose verifier tree has depth $t$ decides on $W$ the language $\{Y=f(X)\}$, or the language $\{Y\le f(X)\}$. Then some nonzero multiple of $Y-f(X)$ in $\RR[X,Y]$ has a straight-line program with at most $4t$ operations. In particular, if the verifier runs in time $t$, this program has $O(t)$ operations.
\end{lemma}

\begin{proof}
\emph{A generic point.} Let $k_0\subseteq\RR$ be the field generated by the constants of the verifier tree and the coefficients of $f$. It is countable. The points of $\RR^n$ whose coordinates are algebraically dependent over $k_0$ form the union of the zero sets of the countably many nonzero polynomials in $k_0[X]$. Each of these zero sets is closed with empty interior, so by the Baire category theorem the union contains no nonempty open set. Fix $x\in U$ outside the union, and put $w_0=(x,f(x))$.

\emph{Paths.} For a coin tape $r\in\{0,1\}^t$ and a sequence of replies $b\in\{0,1\}^t$, let $T_{r,b}$ be the tree in the sense of \cref{def:tree} obtained from the verifier tree by continuing with the branch $r_j$ at the $j$-th coin node of every path, continuing with the branch $b_j$ at the $j$-th receive node of every path, and contracting the send nodes. Its depth is at most $t$, and each of its paths still carries a transcript. Let $\pi_{r,b}$ be the path of $T_{r,b}$ that $w_0$ follows, and represent the values computed on it as fractions $P/D$ of polynomials in $(X,Y)$ as in \cref{sec:cells}. At a test node of $\pi_{r,b}$ call the numerator $P$ of the tested value the \emph{tested numerator}. Since $w_0$ follows $\pi_{r,b}$, every denominator on it is nonzero at $w_0$. So the polynomial $S=PD$ of condition \textup{(b)} of \cref{sec:cells} vanishes at $w_0$ exactly when $P$ does, and it is the zero polynomial exactly when $P$ is.

\emph{Some tested numerator vanishes nontrivially at $w_0$.} We claim that for some pair $(r,b)$ some tested numerator on $\pi_{r,b}$ is a nonzero polynomial that vanishes at $w_0$. Suppose not, and consider one of the finitely many pairs $(r,b)$. At every division on $\pi_{r,b}$ the numerator $P_2$ of the divisor is nonzero at $w_0$. At every test on $\pi_{r,b}$, either the tested numerator is nonzero at $w_0$, and then so is $S$, or the tested numerator is the zero polynomial, and then so is $S$. A polynomial that is nonzero at $w_0$ keeps its sign on a neighbourhood of $w_0$, and the zero polynomial vanishes everywhere. By conditions \textup{(a)}--\textup{(c)} of \cref{sec:cells}, every point of some neighbourhood of $w_0$ follows $\pi_{r,b}$ in $T_{r,b}$. Let $V_0$ be the intersection of these neighbourhoods over all pairs $(r,b)$. Inside $V_0$, for each coin tape and each sequence of replies, the verifier sends the same bits and reaches the same verdict as at $w_0$.

\emph{Freezing the honest prover.} Since $w_0\in W$ belongs to the language, some prover $\mathcal P$ has acceptance probability at least $2/3$ at $w_0$. Let $\mathcal P'(w,s)=\mathcal P(w_0,s)$ for every input $w$ and every transcript $s$. This prover ignores its input and plays the table of replies of $\mathcal P$ at $w_0$, which is a strategy on the finitely many transcripts of this input size; it does not need to know the coin tape. Let $w\in V_0$ and $r\in\{0,1\}^t$. Let $b\in\{0,1\}^t$ consist of the replies of $\mathcal P$ in the interaction at $w_0$ with coin tape $r$, followed by zeros. That interaction follows the path of the verifier tree that corresponds to $\pi_{r,b}$. By induction along this path, the interaction of $\mathcal P'$ at $w$ with coin tape $r$ follows the same path: at computation and test nodes because $w$ follows $\pi_{r,b}$ in $T_{r,b}$, at coin nodes because the coin tape is the same, and at a receive node because the transcript up to that node is the same as at $w_0$, so that $\mathcal P'$ sends the bit that $\mathcal P$ sent at $w_0$. Hence the interaction ends at the same leaf as at $w_0$, and the acceptance probability of $\mathcal P'$ at $w$ equals that of $\mathcal P$ at $w_0$, which is at least $2/3$. The set $V_0\cap W$ is a neighbourhood of $w_0=(x,f(x))$, so it contains a point $w=(x,y')$ with $y'>f(x)$. This point lies outside both languages, so every prover has acceptance probability at most $1/3$ at $w$. This contradiction proves the claim.

\emph{Division by the graph equation.} Let $P\in\RR[X,Y]$ be a nonzero tested numerator with $P(x,f(x))=0$. It is formed from the input coordinates and the constants of the verifier tree by additions, subtractions and multiplications, so its coefficients lie in $k_0$, and so do those of $g=P(X,f(X))\in k_0[X]$. Since $g(x)=0$ and the coordinates of $x$ are algebraically independent over $k_0$, the polynomial $g$ is zero. Division by $Y-f(X)$, which is monic in $Y$, gives $P=(Y-f(X))G+R$ in $\RR[X][Y]$ with $R\in\RR[X]$, and the substitution $Y=f(X)$ gives $R=g=0$. So $P=(Y-f(X))\,G$, and $G\ne0$ because $P\ne0$. By \cref{sec:cells}, the numerators and denominators along $\pi_{r,b}$ up to the test are computed by one straight-line program with at most $4$ operations per computation node, hence with at most $4t$ operations, whose constants are constants of the verifier tree. This program computes $P$. A verifier running in time $t$ unrolls into a verifier tree of depth $O(t)$, which gives the last assertion.
\end{proof}

The number of pairs $(r,b)$ may be exponential in $t$. This costs nothing in the bound, because only one numerator is extracted. The restriction to Boolean messages is used at the step that makes the paths locally constant: if the honest prover sends real numbers that vary with $x$, its possible messages already form an infinite set at a fixed input size.

\begin{theorem}\label{thm:bip}
For every $m$ let a Boolean-message protocol whose verifier tree has depth $t_m$ decide $\{Y=\per_m(X)\}$ on $\RR^{m^2}\times\RR$. Then $t_m\ge2^{\Omega((\log m)^{4/3})}$. Hence every Boolean-message protocol for $\{Y=\per_m(X)\}$ has verifier time $2^{\Omega((\log m)^{4/3})}$, the language $\{Y=\per_m(X)\}$ is not in $\BIPR$, and
\[
\BIPR\subsetneq\PARR .
\]
\end{theorem}

The theorem allows private coins, any number of rounds, a computationally unbounded prover and arbitrary real constants in the verifier. By \cref{lem:boolean} it suffices that the protocol decide the language on a neighbourhood of the graph of $\per_m$ over a nonempty open set.

\begin{proof}
Apply \cref{lem:boolean} with $n=m^2$, $f=\per_m$, $U=\RR^{m^2}$ and $W=\RR^{m^2}\times\RR$. Some nonzero multiple of $Y-\per_m(X)$ has a straight-line program with at most $4t_m$ operations. The permanent has degree $m\ge1$ in $m^2$ variables, so by \cref{thm:burgisser} its border size is at most $Cm^3(4t_m+m^2)$. By \cref{thm:lowerbound} its border size exceeds $2^{c(\log m)^{4/3}}$ for $m\ge m_0$. Since $m^{O(1)}=2^{O(\log m)}$, the two bounds together give $t_m\ge2^{\Omega((\log m)^{4/3})}$. A verifier running in time $T$ unrolls into verifier trees of depth $O(T)$, so its running time is $2^{\Omega((\log m)^{4/3})}$ as well. In particular no verifier running in polynomial time decides $\{Y=\per_m(X)\}$, and the language is not in $\BIPR$. It is in $\PARR$ by \cref{lem:ryser}. Together with the inclusion $\BIPR\subseteq\PARR$ of \cite{IdR99}, this proves $\BIPR\subsetneq\PARR$.
\end{proof}

\begin{remark}[other discrete resources]\label{rem:discrete}
\begin{enumerate}[label=\textup{(\alph*)}]
\item The proof of \cref{lem:boolean} applies when the prover's messages are taken from a fixed countable set $\mathcal M\subseteq\RR$ instead of $\{0,1\}$, and the verifier's messages are still bits. The generic point $x$ can then be chosen with coordinates algebraically independent over the field generated by the constants of the verifier, the coefficients of $f$ and the elements of $\mathcal M$, which is still countable. There are now infinitely many sequences of replies, but only the finitely many paths followed at $w_0$ under the frozen strategy, one for each coin tape, enter the argument; their tested numerators have coefficients in that field, and their programs use the received messages as constants.
\item The proof also covers any purely discrete resource, for instance Boolean nondeterminism, Boolean alternation, or several provers with Boolean messages. The argument uses only that, unless some tested numerator vanishes nontrivially at the generic point $w_0$, the finitely many paths given by the choices of the discrete data are followed on a whole neighbourhood of $w_0$. The outcome of each choice is then locally constant, and so is every acceptance condition formed from these outcomes, including the acceptance probability of frozen strategies; but $w_0$ is accepted and every neighbourhood of $w_0$ contains rejected inputs.
\end{enumerate}
\end{remark}

\subsection{Constant-round Arthur--Merlin proofs}

\begin{definition}\label{def:am}
Let $c,\rho,\mu\ge1$. A \emph{$c$-round public-coin protocol} on inputs $x\in\RR^n$ runs as follows. In round $i=1,\dots,c$ Arthur sends a string $r_i\in\{0,1\}^\rho$ of fair coin flips and Merlin answers with $z_i\in\RR^\mu$. Then a tree $V$ of depth $t$, the \emph{verifier}, accepts or rejects $(x,r_1,z_1,\dots,r_c,z_c)$. For a history $h=(r_1,z_1,\dots,r_{i-1},z_{i-1})$ of the first $i-1$ rounds put $\val_{c+1}(x,h)=1$ if $V$ accepts $(x,h)$ and $\val_{c+1}(x,h)=0$ otherwise, and
\begin{equation}\label{eq:amval}
\val_i(x,h)=\EE_{r_i}\ \sup_{z_i\in\RR^\mu}\ \val_{i+1}(x,h,r_i,z_i)\qquad(i=c,\dots,1),
\end{equation}
where $r_i$ is uniform in $\{0,1\}^\rho$. The protocol \emph{decides} $L\subseteq\RR^n$ if $\val_1(x)\ge2/3$ for $x\in L$ and $\val_1(x)\le1/3$ for $x\notin L$.
\end{definition}

By downward induction on $i$, every $\val_i$ takes values in the finite set $2^{-\rho(c+1-i)}\ZZ\cap[0,1]$. So the suprema are attained, and $\val_1(x)$ is the acceptance probability against the best Merlin. The coins are public: Merlin's answer $z_i$ may depend on $x$ and on everything sent before it. A protocol in which Merlin speaks first is the case in which $V$ ignores $r_1$.

A language $L\subseteq\RR^\infty$ is in $\AMR[c]$ if there are a machine running in polynomial time and polynomial bounds on $\rho$ and $\mu$ in terms of the input size $n$ such that, for every $n$, the $c$-round protocol whose verifier $V$ is the unrolling of the machine on inputs $(x,r_1,z_1,\dots,r_c,z_c)$ with $x\in\RR^n$ decides $L\cap\RR^n$. Put $\AMR[O(1)]=\bigcup_c\AMR[c]$, and let $\AMR[\mathrm{poly}]$ be the class defined in the same way with a number of rounds bounded by a polynomial in $n$. We write $\IPR$ for the class of real interactive proofs of Baartse and Meer \cite{BM16}.

\begin{theorem}\label{thm:am}
For every constant $c\ge1$,
\[
\AMR[c]=\NPR .
\]
Hence
\[
\AMR[O(1)]=\NPR\subsetneq\AMR[\mathrm{poly}]=\IPR .
\]
\end{theorem}

The equality $\AMR[c]=\NPR$ is proved below by fixing the coins, and its proof uses no lower bound; in particular it does not use \cref{thm:lowerbound}. The strict inclusion and the last equality are due to Baartse and Meer, building on Ivanov and de~Rougemont \cite{IdR99}: by \cite[Theorem~14]{BM16} we have $\PARR\subsetneq\IPR$, so $\NPR\subseteq\PARR\subsetneq\IPR$, and by \cite[Proposition~6]{BM16} public coins suffice for $\IPR$, so $\AMR[\mathrm{poly}]=\IPR$. Thus a constant number of rounds of public coins adds nothing to a real certificate, while polynomially many rounds lead beyond $\PARR$.

\begin{lemma}[fixing the coins]\label{lem:coins}
Let a $c$-round public-coin protocol on inputs in $\RR^n$ be given, with $\rho$ coins per round, messages in $\RR^\mu$ and a verifier tree $V$ of depth $t$. There are multisets $R_1,\dots,R_c$ of elements of $\{0,1\}^\rho$, of sizes $q_1,\dots,q_c\ge2$ with
\[
q_1\cdots q_c\le A^{3^c},\qquad A=O\bigl(c^2(n+c\mu)\,\mu\,(\rho c+t)\bigr),
\]
where the constant in $O$ is absolute, such that the \emph{sampled values}
\begin{equation}\label{eq:amsampled}
\val'_{c+1}=\val_{c+1},\qquad
\val'_i(x,h)=\frac1{q_i}\sum_{r_i\in R_i}\ \sup_{z_i\in\RR^\mu}\ \val'_{i+1}(x,h,r_i,z_i)\qquad(i=c,\dots,1),
\end{equation}
with the sum counting multiplicities, satisfy $|\val'_1(x)-\val_1(x)|\le1/8$ for all $x\in\RR^n$.
\end{lemma}

\begin{proof}
\emph{Construction.} Put $\gamma=1/(8c)$. We choose the multisets from the last round to the first. Suppose that $R_{i+1},\dots,R_c$ have been chosen, and put $P_{i+1}=q_{i+1}\cdots q_c$, with $P_{c+1}=1$. Then $\val'_{i+1}$ is defined and takes values in $P_{i+1}^{-1}\ZZ\cap[0,1]$, so suprema of it are attained. As $(x,h)$ ranges over all inputs and all histories of the first $i-1$ rounds, real messages included, the functions
\[
g_{x,h}\colon\{0,1\}^\rho\to[0,1],\qquad g_{x,h}(r_i)=\sup_{z_i\in\RR^\mu}\val'_{i+1}(x,h,r_i,z_i),
\]
form a family $\mathfrak F_i$. It is finite, because each of its members maps a finite set into a finite set, and it does not depend on $R_i$. Let $q_i$ be the least integer with $q_i>32c^2\ln(2|\mathfrak F_i|)$; in particular $q_i\ge2$. Draw $q_i$ independent uniform elements of $\{0,1\}^\rho$. For each $g\in\mathfrak F_i$, by Hoeffding's inequality \cite{Hoe63} the sample mean of $g$ differs from $\EE g$ by at least $\gamma$ with probability at most $2\exp(-2q_i\gamma^2)=2\exp(-q_i/(32c^2))<1/|\mathfrak F_i|$. By the union bound some outcome $R_i$ has a sample mean within $\gamma$ of the mean for every $g\in\mathfrak F_i$. Fix such an $R_i$.

\emph{Accumulated error.} We show $|\val'_i(x,h)-\val_i(x,h)|\le(c+1-i)\gamma$ for all $x$ and $h$, by downward induction on $i$. For $i=c+1$ the two sides agree. Let $g(r_i)=\sup_{z_i}\val_{i+1}(x,h,r_i,z_i)$ and $g'=g_{x,h}$. Suprema do not increase sup-norm distances, so $|g'-g|\le(c-i)\gamma$ pointwise by the induction hypothesis, and averages do not increase them either. Hence
\[
|\val'_i(x,h)-\val_i(x,h)|\le\Bigl|\frac1{q_i}\sum_{r_i\in R_i}g'(r_i)-\EE g'\Bigr|+\bigl|\EE g'-\EE g\bigr|\le\gamma+(c-i)\gamma .
\]
For $i=1$ this gives $|\val'_1-\val_1|\le c\gamma=1/8$.

\emph{Counting.} It remains to bound $|\mathfrak F_i|$. The suprema in \eqref{eq:amsampled} are attained and range over independent arguments, so a sum of suprema is a maximum over joint choices. Hence, for $\theta\in P_{i+1}^{-1}\ZZ$, the statement $g_{x,h}(r_i)\ge\theta$ says: there are a reply $z_i$ and replies for every sampled continuation, that is, a reply $z_{i+1}$ for every element of $R_{i+1}$, a reply $z_{i+2}$ for every pair in $R_{i+1}\times R_{i+2}$, and so on, such that at least a fraction $\theta$ of the $P_{i+1}$ sampled leaves $(x,h,r_i,z_i,r_{i+1},z_{i+1},\dots,r_c,z_c)$ with $(r_{i+1},\dots,r_c)\in R_{i+1}\times\dots\times R_c$ are accepted by $V$. Since every $q_j\ge2$, this is one block of at most $K=2\mu P_{i+1}$ existentially quantified reals. Whether $V$ accepts an input is a Boolean combination of sign conditions of the polynomials of conditions \textup{(a)}--\textup{(c)} of \cref{sec:cells}, at most one for each node of $V$, and $V$ has fewer than $2^t$ nodes. Along a path each operation at most doubles the degrees of the numerators and denominators, which therefore stay at most $2^t$, so these polynomials, including the products $S=PD$, have degree at most $\Delta=2^{t+1}$. Substitute the coin bits: there are $2^{\rho i}$ coin prefixes $(r_1,\dots,r_i)$ and $P_{i+1}$ sampled continuations. So, over all coin prefixes and all $r_i$, the statements involve at most $S=2^{\rho i}P_{i+1}2^{t}$ polynomials of degree at most $\Delta$ in the $M=n+(i-1)\mu\le n+c\mu$ free reals $(x,z_1,\dots,z_{i-1})$ and the $K$ quantified ones. We use the following counting bound.
\begin{enumerate}[label=\textup{(C)},leftmargin=3em]
\item For $S$ polynomials of degree at most $\Delta$ in $M+K$ variables, $\RR^M$ is partitioned into at most $(S\Delta)^{O(MK)}$ sets on each of which the truth value of every one-block existential statement over these polynomials is constant.
\end{enumerate}
Here a one-block existential statement over the polynomials is a formula $\exists u\in\RR^K\,\Psi$, where $\Psi$ is a Boolean combination of the conditions that a polynomial is positive, zero or negative at $(y,u)$, with $y\in\RR^M$ free. The bound \textup{(C)} follows from block elimination \cite[Algorithm~14.1]{BPR06}. It produces a family of $S^{K+1}\Delta^{O(K)}$ polynomials in the $M$ free variables, of degree $\Delta^{O(K)}$, such that the set of sign conditions realised by the original polynomials on the fiber $\{y\}\times\RR^K$ is constant as $y$ ranges over a connected component of a sign condition of the family. By the usual count of sign conditions \cite[Proposition~7.31]{BPR06}, the number of these components is at most $(S\Delta)^{O(MK)}$. The set of sign conditions realised on the fiber determines the truth value of every one-block existential statement.

The value $g_{x,h}(r_i)$ is the largest $\theta$ for which $g_{x,h}(r_i)\ge\theta$ holds. So $g_{x,h}$ is determined by the coin prefix of $h$ and by the set of the partition in \textup{(C)} that contains $(x,z_1,\dots,z_{i-1})$. Hence $|\mathfrak F_i|\le2^{\rho c}(S\Delta)^{O(MK)}$ and
\[
\ln|\mathfrak F_i|\le O(MK)\,(\rho c+t+\log P_{i+1}).
\]
Since $q_i\le32c^2\ln(2|\mathfrak F_i|)+1$, $K\le2\mu P_{i+1}$ and $M\le n+c\mu$, this gives $q_i\le A\,P_{i+1}(1+\ln P_{i+1})\le A\,P_{i+1}^2$ for a suitable $A=O\bigl(c^2(n+c\mu)\mu(\rho c+t)\bigr)$. Then $P_i=q_iP_{i+1}\le A\,P_{i+1}^3$, so $P_i\le A^{(3^{c+1-i}-1)/2}$ by downward induction, and $q_1\cdots q_c=P_1\le A^{3^c}$.
\end{proof}

Write each multiset as a list $R_i=(r_i^{(1)},\dots,r_i^{(q_i)})$ and put $P_1=q_1\cdots q_c$. A \emph{sampled leaf} is a tuple $(j_1,\dots,j_c)$ with $1\le j_i\le q_i$; there are $P_1$ of them. A \emph{strategy tree} assigns a reply $z_1(j_1)\in\RR^\mu$ to each element of $R_1$, a reply $z_2(j_1,j_2)\in\RR^\mu$ to each pair in $R_1\times R_2$, and so on up to $z_c(j_1,\dots,j_c)$. It is a vector of $\mu(q_1+q_1q_2+\dots+P_1)\le2\mu P_1$ reals. The sampled leaf $(j_1,\dots,j_c)$ \emph{accepts} if $V$ accepts $\bigl(x,r_1^{(j_1)},z_1(j_1),\dots,r_c^{(j_c)},z_c(j_1,\dots,j_c)\bigr)$. Since the suprema in \eqref{eq:amsampled} range over independent subtrees and are attained, downward induction on the round gives
\begin{equation}\label{eq:strategytree}
\val'_1(x)=\max_{\text{strategy trees}}\ \frac{\#\{\text{accepting sampled leaves}\}}{P_1} .
\end{equation}
So $\val'_1(x)\ge1/2$ holds exactly when some strategy tree makes at least half of the $P_1$ sampled leaves accept. Let $T$ be the tree on inputs $(x,z)$, with $z$ a strategy tree, that runs $V$ on every sampled leaf in turn, with the coin bits of the leaf as constants, counts the accepting runs, and accepts when their number is at least $P_1/2$. Before each division in a simulated run it tests the divisor and counts that leaf as rejecting if the divisor is zero; by the convention of \cref{sec:trees} these tests are already present in $V$. The depth of $T$ is $O(P_1t)$. If the protocol decides $L$, then by \cref{lem:coins} we have $\val'_1\ge2/3-1/8>1/2$ on $L$ and $\val'_1\le1/3+1/8<1/2$ off $L$, so
\begin{equation}\label{eq:amcert}
T\text{ accepts }(x,z)\text{ for some strategy tree }z\quad\Longleftrightarrow\quad x\in L .
\end{equation}

\begin{proof}[Proof of \cref{thm:am}]
Let $L\in\AMR[c]$ with $c$ constant. The numbers $\rho$ and $\mu$ and the depth $t$ of the verifier tree are polynomial in the input size $n$, so $A$ is polynomial in $n$, and so is $P_1\le A^{3^c}$. The multisets of \cref{lem:coins} for all input sizes form Boolean advice of polynomial length. A machine can carry all of it in one real constant and read the bits by doubling and comparing, as in the proof that randomised polynomial time over the reals equals $\PR$ \cite{CKKLW95,Koi94}. A machine that reads the multisets for the size of its input, takes a strategy tree $z$ of at most $2\mu P_1$ reals as its certificate, and runs $T$ on $(x,z)$ works in polynomial time, and by \eqref{eq:amcert} it witnesses $L\in\NPR$. Conversely, every language in $\NPR$ is in $\AMR[c]$: Merlin sends a certificate in the first round, and the verifier ignores the coins and the other rounds.
\end{proof}

\begin{theorem}\label{thm:amper}
For every $m$ let a public-coin protocol with $c_m$ rounds, real messages from Merlin and a verifier running in time polynomial in $m$ decide $\{Y=\per_m(X)\}$ on $\RR^{m^2}\times\RR$. Then $c_m\ge\Omega(\log\log m)$. In particular $\{Y=\per_m(X)\}$ is not in $\AMR[O(1)]$.
\end{theorem}

\begin{proof}
Fix $m$, write $c=c_m$, and let $\rho$, $\mu$ and $t$ be the number of coins per round, the length of Merlin's messages and the depth of the verifier tree, all polynomial in $m$. If $c>m$ there is nothing to prove, so let $c\le m$. Let $P_1$ and $T$ be as above, for the input size $m^2+1$. By \eqref{eq:amcert}, the tree $T$ satisfies the hypothesis of \cref{thm:graph} with $f=\per_m$, $n=m^2$, $d=m$, $U=\RR^{m^2}$, depth $O(P_1t)$ and at most $2\mu P_1$ certificate coordinates. So the border size of $\per_m$ is bounded by a polynomial in $m$, $t$ and $P_1$, and hence, as $t$ is polynomial in $m$, by $m^{O(1)}P_1^{O(1)}$. By \cref{thm:lowerbound} it exceeds $2^{\Omega((\log m)^{4/3})}$ for large $m$, which forces $P_1\ge2^{\Omega((\log m)^{4/3})}$. Since $n=m^2+1$ and $c$, $\rho$, $\mu$ and $t$ are polynomial in $m$, we have $\log A=O(\log m)$. From $P_1\le A^{3^c}$ we get $3^c\ge\Omega((\log m)^{1/3})$, and so $c\ge\Omega(\log\log m)$.
\end{proof}

The protocol of Lund, Fortnow, Karloff and Nisan \cite{LFKN92} verifies the real permanent with public coins, real messages and $m$ rounds. The lower bound of \cref{thm:amper} is doubly logarithmic because the sample sizes in \cref{lem:coins} grow by a fixed power at each round. \Cref{thm:am,thm:amper} concern public coins; for Boolean messages, \cref{thm:bip} excludes protocols with private coins and any number of rounds.

\section{Open problems}\label{sec:open}

\paragraph{$\PR$ versus $\NPR$.} If a language with short real certificates is the graph of a polynomial of polynomial degree, or the region on one side of such a graph, then \cref{thm:graph,thm:threshold} give an upper bound on the border complexity of that polynomial. Therefore a lower bound on border complexity cannot prove that such a language lies outside $\PR$. If a tree of depth $t$ decides a language whose boundary contains an irreducible hypersurface $\{h=0\}$, then the generic-path argument gives a nonzero multiple of $h$ that has a straight-line program of length $O(t)$, but with no bound on its degree \cite{Bur04}.

B\"urgisser conjectures that $\VNP\not\subseteq\overline{\VP}$ over $\CC$ implies $\PC\ne\NPC$ nonuniformly \cite[Conjecture~1.1]{Bur26}. The hypothesis follows from \cref{thm:lowerbound} in its form over $\CC$. Polynomial-size circuits for the resultant of $n+1$ quadratic forms imply $\VP=\VNP$ over $\CC$ \cite{Bur26}. Extending this implication to every nonzero multiple of that resultant would prove the conjecture. The resultant has exponential degree, so the degree dependence in \cref{thm:burgisser} is an obstacle. See also \cite{SS95} for the analogous role of multiples of integers.

\paragraph{B\"urgisser's conjecture and squarefree parts.} We record a side observation: B\"urgisser's conjecture follows from a bound on the circuit complexity of squarefree parts. Unlike the rest of the paper, this observation is not formalized in Lean. For a nonzero polynomial $P$ over $\CC$ write $\rad(P)$ for the product of its distinct irreducible factors, and write $\doteq$ for equality up to a nonzero scalar. Consider the following hypothesis.
\begin{enumerate}[label=\textup{(R)},leftmargin=3em]
\item There is a polynomial $\pi$ such that, if a polynomial $P$ over $\CC$ is computed by an arithmetic circuit of size $s$, then $\rad(P)$ is computed by an arithmetic circuit of size $\pi(s)$, whatever the degree of $P$.
\end{enumerate}
Let $\overline{\mathrm{(R)}}$ be the weaker hypothesis that bounds the border complexity of $\rad(P)$ by $\pi(s)$. Dutta, Saxena and Sinhababu conjecture a polynomial bound for the complexity of the squarefree part \cite[Conjecture~1]{DSS22}. In a form that bounds the minimum of the degree and the circuit size of $\rad(P)$, it implies \textup{(R)}: the polynomial $P\cdot(Z^{2^s}-1)$, with a new variable $Z$, has a circuit of size $O(s)$ and a squarefree part of degree at least $2^s$, so for large $s$ the bound applies to its circuit size, and setting $Z=0$ gives $\rad(P)$ up to a scalar.

\begin{lemma}\label{lem:radpath}
Let $f\in\CC[x]$ be irreducible, and let a complex tree of depth $t$ decide $\{f=0\}$. Then there are nonzero polynomials $A,B,C\in\CC[x]$, computed together by one straight-line program with $O(t)$ operations, such that
\[
f\ \doteq\ \frac{\rad(AB)\,\rad(CB)}{\rad(ACB)\,\rad(B)} .
\]
\end{lemma}

\begin{proof}
Run the tree at the generic point of $V(f)$, that is, over the function field of the irreducible hypersurface $V(f)$, and represent the values on the path it follows by fractions as in \cref{sec:cells}. The points of a dense open subset of $V(f)$ follow the same path, so the path accepts. Its conditions are $E_1=\dots=E_m=0$ and $B\ne0$, where the $E_i$ are the tested numerators of the tests answered ``yes'' and $B$ is the product of the other tested numerators and of the numerators of the divisors. All of them are computed by one straight-line program with $O(t)$ operations, $f$ divides every $E_i$, and $f$ does not divide $B$. Every point that satisfies the conditions is accepted, so it lies in $V(f)$. We discard the $E_i$ that are the zero polynomial. Some $E_i$ remains, since otherwise the open set $\{B\ne0\}$ would be accepted.

Let $A$ be one of the remaining $E_i$. Let $C$ be a nonzero linear combination of the $E_i$ that is not divisible by any irreducible factor of $A$ that fails to divide some $E_j$. Each such factor excludes a proper linear subspace of coefficient vectors, and there are finitely many. The common irreducible factors of $A$ and $C$ are then exactly the common irreducible factors $g$ of all the $E_i$. If such a $g$ does not divide $B$, then a dense open subset of $V(g)$ satisfies the conditions of the path, so $V(g)\subseteq V(f)$ and $g\doteq f$. Conversely $f$ divides $A$ and $C$ and does not divide $B$. For an irreducible $g$ the exponent of $g$ in the displayed quotient is
\[
\mathbf 1_{g\mid AB}+\mathbf 1_{g\mid CB}-\mathbf 1_{g\mid ACB}-\mathbf 1_{g\mid B},
\]
which is $1$ if $g$ divides $A$ and $C$ but not $B$, and $0$ otherwise.
\end{proof}

\begin{proposition}\label{prop:radical}
Assume \textup{(R)}. If $\NPC$ is contained in nonuniform polynomial time over $\CC$, then $\VP=\VNP$ over $\CC$. Under $\overline{\mathrm{(R)}}$ the same assumption gives $\VNP\subseteq\overline{\VP}$ over $\CC$. In particular $\overline{\mathrm{(R)}}$ implies B\"urgisser's conjecture \cite[Conjecture~1.1]{Bur26}, and together with \cref{thm:lowerbound} in its form over $\CC$ it implies $\PC\ne\NPC$ nonuniformly.
\end{proposition}

\begin{proof}
Let $p(T')$ be a monic polynomial of degree $n$ with roots $r_1,\dots,r_n$, let $q_1,\dots,q_n$ have degree less than $n$, and let $T$ be a further variable. Put
\[
f(p,q,T)=\prod_{\sigma\in S_n}\Bigl(1-T\prod_{i=1}^nq_i(r_{\sigma(i)})\Bigr).
\]
Its coefficients are symmetric in the $r_j$, so $f$ is a polynomial in $T$ and in the coefficients of $p$ and of the $q_i$. Treat the $r_j$ and the coefficients of the $q_i$ as independent variables. The values $q_i(r_j)$ are then independent over $\CC(r)$ by the Vandermonde determinant, so the $n!$ products $\prod_iq_i(r_{\sigma(i)})$ are distinct, and $S_n$ permutes them transitively. Their common minimal polynomial over the field generated by the coefficients of $p$ and of the $q_i$ therefore has degree $n!$. It equals $T^{n!}f(1/T)=\prod_\sigma\bigl(T-\prod_iq_i(r_{\sigma(i)})\bigr)$, which is monic in $T$ and hence irreducible as a polynomial by Gauss's lemma. Since $f(0)=1$, also $f$ is irreducible. The set $\{f=0\}$ is in $\NPC$: guess $z_1,\dots,z_n$, check the coefficient identity $p(T')=\prod_j(T'-z_j)$, and check $T\prod_iq_i(z_i)=1$.

By assumption, complex trees of depth polynomial in $n$ decide $\{f=0\}$. \Cref{lem:radpath} and \textup{(R)} give polynomial-size circuits for $U=\rad(AB)\rad(CB)$ and $V=\rad(ACB)\rad(B)$, with $U=\lambda fV$ for some $\lambda\ne0$. Each radical is divisible at most once by $T$, so $V$ has order $a\le2$ in $T$ at $0$. Write $U_j$ and $V_j$ for the coefficients of $T^j$. Since $f(0)=1$ we have $U_a=\lambda V_a\ne0$ and
\[
[T]\,f=\frac{U_{a+1}}{U_a}-\frac{V_{a+1}}{V_a}.
\]
Fix distinct complex numbers $r_1,\dots,r_n$ for which $V_a$ is a nonzero polynomial in the coefficients of the $q_i$, and determine the $q_i$ by interpolation from $q_i(r_j)=X_{ij}$. Then $[T]f=-\sum_\sigma\prod_iX_{i\sigma(i)}=-\per_n(X)$. The coefficients of $T^j$ for $j\le3$ are computed from the circuits for $U$ and $V$ in $\CC[x][T]/(T^4)$. Choose $X_0$ with $V_a(X_0)\ne0$ and evaluate the right-hand side at $X_0+\tau W$ modulo $\tau^{n+1}$; the constant terms of the denominators are nonzero constants, so the divisions cost $O(n^2)$ operations each. Summing the coefficients gives $-\per_n(X_0+W)$ exactly, a polynomial-size circuit for the permanent. As the permanent is complete for $\VNP$, $\VP=\VNP$. Under $\overline{\mathrm{(R)}}$ the same construction applied to approximating circuits gives coefficients that converge to those of the permanent, because the denominators have nonzero limits at $X_0$. Hence the permanent has polynomial border complexity and $\VNP\subseteq\overline{\VP}$.
\end{proof}

The identity of \cref{lem:radpath} holds unconditionally. Hypothesis \textup{(R)} is used only to bound the coefficients of $T^j$, $j\le3$, of the four squarefree parts on the interpolation slice. The identity depends on discarding multiplicities: with the polynomials $AB$, $CB$, $ACB$, $B$ in place of their squarefree parts the quotient is $1$, and quantities that are additive under products, such as root sums counted with multiplicity, cancel in the same way.

Lower bounds for all nonzero elements are known for other ideals. Forbes, Shpilka, Tzameret and Wigderson \cite{FSTW16} prove lower bounds for all nonzero multiples of explicit polynomials in several restricted circuit models. Andrews and Forbes \cite{AF22} show that every nonzero polynomial in the ideal of $r\times r$ minors of a generic matrix yields, through a small oracle circuit of depth three, an approximation of a determinant of size $\Theta(r^{1/3})$. For the principal ideal $(Y-\per_m(X))$ over $\RR$, \cref{thm:lowerbound,thm:burgisser} give a lower bound even for unrestricted arithmetic circuits. If a nonzero multiple has a straight-line program of length $\ell$, then the border complexity of $\per_m$ is at most $Cm^3(\ell+m^2)$. Hence $\ell\ge 2^{\Omega((\log m)^{4/3})}$ regardless of the degree of the multiple.

\paragraph{The universal--existential level.} Suppose $y=f(x)$ holds exactly when for all $z$ some $u$ is accepted. Off the graph there is a $z$ for which every $u$ is rejected. Such witnesses can escape to infinity as $(x,y)$ approaches the graph. For example, let the verifier test $yz-1=0$, accepting exactly when this test fails and ignoring $u$. Then $\forall z\,\exists u$ defines $y=0$. For $y\ne0$ the only rejecting value is $z=1/y$. On any fixed closed ball $B$ in $z$, the accepting margin of \eqref{eq:margin} is $v(y,z)=\min\{1,(1-yz)^2\}$, which tends uniformly to $1$ as $y\to0$. For fixed $N$ and $a$, the objective \eqref{eq:J} at maximizing margins changes from its value at $y=0$ by at most $\sup_{z\in B}|v(y,z)^N-1|$, uniformly over all atom locations and weights in that ball. Its maximum is therefore continuous across $y=0$. It remains open whether the border complexity of $f$ is polynomially bounded in its degree, the number of variables, the tree depth and the two block lengths for this quantifier order.

\paragraph{Thresholds at the second level.} For the permanent, negating one row reduces either closed threshold to the other, and \cref{cor:second} gives the lower bound for both. For a general polynomial $f$, does an $\exists z\,\forall u$ verifier for $\{y\le f(x)\}$ yield the same polynomial bound on border complexity as in \cref{thm:second}? Here the pairs $(y,z)$ over $x$ for which every $u$ is accepted can form a positive-dimensional set. If $v^N$ is bounded by a multiple of the squared distance to this set, the envelope vanishes exactly on it. One possible approach is to minimize $e-\vartheta y$ with a smaller generic linear perturbation to select a boundary point.

\paragraph{Higher levels.} At the third level, the corresponding set of lower dimension is the set of points at which an accepting margin is positive for every value of an intermediate block of variables. Its complement is not described by a single maximization. It is open whether the graph of the permanent lies in $\PHR$. If it does not, then $\PHR\subsetneq\PARR$. It is also open whether $\PHR$ is contained in polynomial time with an oracle for the permanent; such an inclusion would be a real analogue of Toda's theorem. Compare \cite{BM16,Bur26}.

\paragraph{Rounds.} \Cref{thm:amper} gives a lower bound of $\Omega(\log\log m)$ on the number of rounds of a public-coin protocol for the permanent, while the upper bound is $m$. The lower bound is limited by the growth of the sample sizes in \cref{lem:coins}. We do not know whether a constant number of rounds with private coins and real messages can be reduced to a real certificate.


\appendix
\section{Formal verification in Lean}\label{sec:lean}
\begingroup\raggedright

The results of this paper are formalized in the proof assistant Lean~4 (version 4.36.0), using the Mathlib library (version 4.36.0). The formalization is in the directory \texttt{lean/} of the source repository; it has about 36\,000 lines. It is built with \texttt{lake build} and contains no \texttt{sorry}. Every theorem below depends only on the three standard axioms of Lean and Mathlib (\texttt{propext}, \texttt{Classical.choice}, \texttt{Quot.sound}), as reported by \texttt{\#print axioms}. All Lean names are in the namespace \texttt{ClassSep}.

\subsection{What is formalized}

Every numbered lemma, proposition, theorem and corollary of \crefrange{sec:prelim}{sec:interaction} is stated and proved in Lean, as are \cref{thm:A,thm:B,thm:C,thm:D,thm:E,thm:F}, \cref{rem:several}, \cref{rem:complexcoef} and part (a) of \cref{rem:discrete}. The informal part (b) of \cref{rem:discrete}, the example after \cref{prop:value}, the remark after \cref{thm:second} on the length of the existential block, and \cref{sec:open} are not formalized.

\subsection{Cited results as hypotheses}

The results that the paper cites from the literature are not proved in Lean. They enter as explicit hypotheses of the theorems that use them, in the form of propositions, and no axiom is added. They are:
\begin{itemize}[leftmargin=1.6em,itemsep=1pt]
\item \texttt{SAFacts}: the facts \ref{S2}--\ref{S8} on semialgebraic sets, each in a precise form (file \texttt{SA/Facts.lean});
\item \texttt{PermLowerBound}: \cref{thm:lowerbound};
\item \texttt{BurgisserThm}: \cref{thm:burgisser};
\item \texttt{EffectiveQE}: effective quantifier elimination \cite[Theorem~14.16]{BPR06}, used in \cref{lem:flatten}(d);
\item \texttt{CountingBound}: the counting bound (C) in the proof of \cref{lem:coins};
\item \texttt{ClassFacts}: the known relations $\NPR,\coNPR,\Sigma_2^{\RR}\subseteq\PARR$, $\NPC,\coNPC\subseteq\PARC$, $\BIPR\subseteq\PARR$ \cite{IdR99}, $\PARR\subsetneq\IPR$ and $\AMR[\mathrm{poly}]=\IPR$ \cite{BM16}.
\end{itemize}
Fact \ref{S1} is not a hypothesis. A set is \emph{defined} to be semialgebraic if it is defined by a first-order formula over $(\RR,+,\times,<)$ with real parameters, so \ref{S1} holds by construction, and each of its applications is a proof that a concrete set has a defining formula. The theorem of Baur and Strassen (\cref{thm:baurstrassen}), Ryser's formula, Hoeffding's inequality (from Mathlib), and the unrolling of machines into trees are proved in Lean. In particular, all lower bounds for the permanent in the formalization are conditional on \texttt{PermLowerBound}, exactly as in the paper.

\subsection{Modelling choices}

\begin{itemize}[leftmargin=1.6em,itemsep=1pt]
\item The dimension of a semialgebraic set is described by the predicate ``some projection to $d$ coordinates has nonempty interior'' (\texttt{DimGE}); \ref{S4}(a), first part, is then a theorem.
\item A tree returns an outcome: the label of the leaf reached, or a division by zero, which rejects. Sets of marked outcomes replace marked leaves, so the convention of \cref{sec:trees} that every rejection happens at a leaf is not needed. Tests may be applied to any operand.
\item Programs and circuits have numbered slots for their constants. A program or circuit \emph{graph} is the object with slots; this expresses statements such as ``the graph depends only on the graph of the program''.
\item Taylor polynomials are defined from iterated Fr\'echet derivatives, and \cref{cor:derivatives} is stated with iterated derivatives, whose components are the partial derivatives $\partial^\alpha$.
\item Machines are real random-access machines whose index registers address the real registers; they simulate the machines of \cite{BCSS98} with polynomial overhead. $\PARR$ and $\PARC$ are defined by uniform families of algebraic circuits of polynomial depth, where a machine computes, in polynomial time, the description of a gate from the binary digits of its index.
\end{itemize}

\subsection{Proofs that differ from the text}

The following proofs in Lean differ from the arguments given in the text; the statements are the same.
\begin{itemize}[leftmargin=1.6em,itemsep=1pt]
\item \cref{lem:isolate} is proved by induction on the number of certificate coordinates, splitting off one coordinate at a time: the projection to $(x,y)$ either contains a box, and the coordinate is fixed, or has finite fibers over an open set, and \ref{S5} applies. Nash manifolds and tangent spaces are not used.
\item In \cref{lem:newton}(b), in the proof of \cref{thm:onesided} and in the proof of \cref{thm:second}, the analytic implicit function theorem is replaced by little-$o$ estimates and the uniqueness of polynomial approximations of order $o(|\xi-\xi_1|^D)$.
\item In \cref{lem:envelope}, the radius $r_0$ is assumed nonnegative explicitly.
\item In the proof of \cref{thm:onesided}, the finiteness of the set $\Delta$ of exceptional levels is replaced by a density argument.
\item In \cref{lem:boolean}, the generic point is chosen outside the zero sets of finitely many nonzero polynomials, which avoids the Baire category theorem; in \cref{rem:discrete}(a) the Baire category theorem is used.
\end{itemize}

\subsection{Correspondence}

For each result we give the Lean name and the file, relative to \texttt{lean/ClassSeparations/}.

\begin{description}[leftmargin=1.6em,itemsep=1pt,font=\normalfont\bfseries]
\item[\Cref{def:tree,def:border}] \texttt{Tree}, \texttt{Tree.run}, \texttt{Tree.depth}; \texttt{IsApprox}, \texttt{BorderLE} (\texttt{Basic/Tree.lean}, \texttt{Basic/Circuits.lean}).
\item[\Cref{thm:baurstrassen}] \texttt{baur\_strassen}, \texttt{baur\_strassen'} (\texttt{Core/BaurStrassen.lean}).
\item[\Cref{lem:cells}] \texttt{cells}, \texttt{cells\_length} (\texttt{Basic/Cells.lean}).
\item[\Cref{lem:derivatives,cor:derivatives}] \texttt{derivatives}, \texttt{derivatives\_cor} (\texttt{Core/Derivatives.lean}).
\item[\Cref{lem:closure}] \texttt{closure\_lemma} (\texttt{Core/Closure.lean}).
\item[\Cref{lem:newton}] \texttt{newton\_circuit}, \texttt{newtonPoly\_eq\_taylorPoly}, \texttt{tendsto\_coeff\_newtonPoly} (\texttt{Core/Newton.lean}).
\item[\Cref{thm:local,lem:minimisers,lem:genericcentre}] \texttt{local\_theorem}, \texttt{minimisers}, \texttt{genericCentre} (\texttt{Core/Local.lean}).
\item[\Cref{lem:isolate}] \texttt{isolate} (\texttt{Cert/Isolate.lean}).
\item[\Cref{thm:graph,rem:several,thm:threshold}] \texttt{graph\_theorem}, \texttt{graph\_theorem\_vec}, \texttt{threshold\_theorem\_le}, \texttt{threshold\_theorem\_ge} (\texttt{Cert/Graph.lean}).
\item[\Cref{prop:value}] \texttt{value\_functions} (\texttt{Cert/Value.lean}).
\item[\Cref{lem:ryser}] \texttt{ryser} (\texttt{Cert/Permanent.lean}); \texttt{permLang\_mem\_PAR}, \texttt{cpermLang\_mem\_PAR} (\texttt{Machines/Perm.lean}).
\item[\Cref{cor:np}] \texttt{cor\_np} (\texttt{Cert/Permanent.lean}) for trees; \texttt{theorem\_A} (\texttt{Main.lean}) for classes.
\item[\Cref{thm:onesided,lem:margin}] \texttt{onesided}, \texttt{Margin.margin\_cert\_iff}, \texttt{OnesidedAux.lemma\_5\_2\_b}, \texttt{OnesidedAux.lemma\_5\_2\_c} (\texttt{Universal/Onesided.lean}, \texttt{Universal/Margin.lean}).
\item[\Cref{thm:complex,rem:complexcoef}] \texttt{complex\_field}, \texttt{complex\_coeff} (\texttt{Universal/Complex.lean}).
\item[\Cref{cor:conp}] \texttt{cor\_conp\_universal}, \texttt{cor\_conp\_strict}, \texttt{cor\_conp\_complex} (\texttt{Universal/Separations.lean}); \texttt{theorem\_B}, \texttt{theorem\_C} (\texttt{Main.lean}).
\item[\Cref{thm:second}] \texttt{second\_level} (\texttt{Second/Main.lean}), from \texttt{sl\_setup} (\texttt{Second/Setup.lean}) and \texttt{sl\_conclude\_proof} (\texttt{Second/Conclude.lean}).
\item[\Cref{lem:goodz,lem:flatten}] \texttt{SetupAux.lemma\_6\_2} (\texttt{Second/Goodz.lean}); \texttt{SetupAux.lemma\_6\_3\_a}, \texttt{\_b}, \texttt{\_c} (\texttt{Second/Flatten.lean}), \texttt{SetupAux.lemma\_6\_3\_d} (\texttt{Second/FlattenD.lean}) with \texttt{exponent\_bound} (\texttt{Second/Exponent.lean}).
\item[\Cref{lem:envelope,lem:continuity,lem:envelopevalue,lem:tilt}] \texttt{envelope\_props} (\texttt{Second/Envelope.lean}), \texttt{SLData.Good.env\_continuity} (\texttt{Second/Continuity.lean}), \texttt{SLData.env\_value} (\texttt{Second/EnvValue.lean}), \texttt{SLData.Good.tilt} (\texttt{Second/Tilt.lean}).
\item[\Cref{cor:second}] \texttt{cor\_second} (\texttt{Second/Main.lean}); \texttt{theorem\_B} (\texttt{Main.lean}).
\item[\Cref{def:bip,lem:boolean,thm:bip}] \texttt{VTree}, \texttt{boundary\_test}, \texttt{bip\_lower} (\texttt{Interaction/Boolean.lean}); \texttt{theorem\_D} (\texttt{Main.lean}).
\item[\Cref{rem:discrete}(a)] \texttt{boundary\_test\_countable} (\texttt{Interaction/Remark74.lean}).
\item[\Cref{def:am,lem:coins,thm:amper}] \texttt{AMProtocol}, \texttt{fixing\_coins}, \texttt{strategy\_tree} (equation \eqref{eq:amcert}), \texttt{am\_rounds} (\texttt{Interaction/AM.lean}).
\item[\Cref{thm:am}] \texttt{am\_eq\_np}, \texttt{am\_const\_eq\_np} (\texttt{Machines/AMNP.lean}).
\item[\Cref{thm:A,thm:B,thm:C,thm:D,thm:E,thm:F}] \texttt{theorem\_A}, \dots, \texttt{theorem\_F} (\texttt{Main.lean}).
\item[Classes] \texttt{P\_R}, \texttt{NP\_R}, \texttt{coNP\_R}, \texttt{Sigma2\_R}, \texttt{PAR\_R} (\texttt{Machines/Classes.lean}); \texttt{NP\_C}, \texttt{coNP\_C}, \texttt{PAR\_C} (\texttt{Machines/Complex.lean}); \texttt{BIP\_R}, \texttt{IP\_R}, \texttt{AM\_R}, \texttt{AM\_R\_const}, \texttt{AM\_R\_poly} (\texttt{Machines/Interactive.lean}).
\end{description}

The file \texttt{Pinned.lean} restates the main theorems verbatim and checks that the proved theorems have exactly these statements.
\par\endgroup


\bibliographystyle{alpha}
\bibliography{refs}
\end{document}
